% ─────────────────────────────────────────────────────────────────────────────
% informe_completo_07_08.tex
% Overleaf: Menu -> Compiler -> pdfLaTeX
% Une Experimento 07 (predicción empírica de moyo) y Experimento 08
% (teoría de invariantes: por qué el campo solo ve 4 de 7 parámetros).
% ─────────────────────────────────────────────────────────────────────────────
\documentclass[11pt, a4paper]{article}

% -- Codificacion --------------------------------------------------------------
\usepackage[T1]{fontenc}
\usepackage[utf8]{inputenc}

% -- Fuentes (disponibles en TeX Live / Overleaf) -------------------------------
\usepackage[default, scale=0.95]{sourcesanspro}
\usepackage[scaled=0.82]{inconsolata}

% -- Geometria -------------------------------------------------------------------
\usepackage[
  a4paper,
  top        = 2.8cm,
  bottom     = 2.4cm,
  left       = 2.2cm,
  right      = 2.2cm,
  headheight = 30pt,
  headsep    = 14pt
]{geometry}

% -- Colores ----------------------------------------------------------------------
\usepackage[table]{xcolor}
\definecolor{cNavy}   {HTML}{0D1B3E}
\definecolor{cBlue}   {HTML}{2A4F9B}
\definecolor{cAccent} {HTML}{4A7FD4}
\definecolor{cMuted}  {HTML}{6B7FA8}
\definecolor{cCodeFG} {HTML}{1B3A70}
\definecolor{cCodeBG} {HTML}{EEF2FA}
\definecolor{cTHbg}   {HTML}{0D1B3E}
\definecolor{cTRalt}  {HTML}{F2F5FC}
\definecolor{cRule}   {HTML}{C0CADF}
\definecolor{cQuoteBG}{HTML}{F0F5FF}
\definecolor{cGood}   {HTML}{1E7A4C}
\definecolor{cBad}    {HTML}{B23A3A}

% -- Matematicas --------------------------------------------------------------------
\usepackage{amsmath}
\usepackage{amssymb}
\usepackage{mathtools}

% -- Tipografia y espaciado -----------------------------------------------------------
\usepackage{microtype}
\usepackage{parskip}
\setlength{\parindent}{0pt}
\setlength{\parskip}{5pt plus 1pt}

% -- Tablas --------------------------------------------------------------------------
\usepackage{booktabs}
\usepackage{tabularx}
\usepackage{colortbl}
\usepackage{array}
\usepackage{adjustbox}
\renewcommand{\arraystretch}{1.4}
\newcolumntype{L}[1]{>{\raggedright\arraybackslash}p{#1}}
\newcolumntype{C}[1]{>{\centering\arraybackslash}p{#1}}
\newcolumntype{R}[1]{>{\raggedleft\arraybackslash}p{#1}}

% -- Cajas de codigo y citas -----------------------------------------------------------
\usepackage{tcolorbox}
\tcbuselibrary{skins, breakable}

\tcbset{
  codebox/.style = {
    enhanced,
    colback    = cCodeBG,
    colframe   = cAccent,
    arc        = 3pt,
    leftrule   = 3pt,  rightrule  = 0.5pt,
    toprule    = 0.5pt, bottomrule = 0.5pt,
    left = 10pt, right = 10pt, top = 7pt, bottom = 7pt,
    fontupper  = \small\ttfamily\color{cCodeFG},
    breakable
  },
  quotebox/.style = {
    enhanced,
    colback    = cQuoteBG,
    colframe   = cAccent,
    arc        = 2pt,
    leftrule   = 3pt,  rightrule  = 0pt,
    toprule    = 0pt,  bottomrule = 0pt,
    left = 12pt, right = 10pt, top = 7pt, bottom = 7pt,
    fontupper  = \itshape\color{cCodeFG},
    breakable
  },
  pendingbox/.style = {
    enhanced,
    colback    = white,
    colframe   = cMuted,
    arc        = 2pt,
    dashed,
    left = 12pt, right = 10pt, top = 7pt, bottom = 7pt,
    fontupper  = \color{cMuted},
    breakable
  },
  defbox/.style = {
    enhanced,
    colback    = white,
    colframe   = cBlue,
    arc        = 2pt,
    leftrule   = 3pt,  rightrule  = 0.5pt,
    toprule    = 0.5pt, bottomrule = 0.5pt,
    left = 12pt, right = 10pt, top = 7pt, bottom = 7pt,
    breakable
  }
}

% -- Encabezados de seccion -------------------------------------------------------------
\usepackage{titlesec}

\titleformat{\part}
  {\color{cNavy}\huge\bfseries}
  {\color{cAccent}Parte \thepart}{0.6em}{}
  [\vspace{2pt}{\color{cAccent}\hrule height 2pt}]

\titleformat{\section}
  {\color{cNavy}\Large\bfseries}
  {}{0em}{}
  [\vspace{2pt}{\color{cAccent}\hrule height 1.6pt}\vspace{4pt}]

\titleformat{\subsection}
  {\color{cBlue}\large\bfseries}
  {}{0em}{}
  [\vspace{1pt}{\color{cRule}\hrule height 0.5pt}\vspace{2pt}]

\titleformat{\subsubsection}
  {\color{cBlue}\normalsize\bfseries}
  {}{0em}{}

\titlespacing*{\section}      {0pt}{20pt}{8pt}
\titlespacing*{\subsection}   {0pt}{14pt}{4pt}
\titlespacing*{\subsubsection}{0pt}{10pt}{3pt}

% -- Cabecera y pie -----------------------------------------------------------------
\usepackage{fancyhdr}
\usepackage{tikz}
\usetikzlibrary{calc}

\pagestyle{fancy}
\fancyhf{}
\renewcommand{\headrulewidth}{0pt}

\fancyhead[C]{%
  \begin{tikzpicture}[remember picture, overlay]
    \fill[cNavy]
      (current page.north west)
      rectangle ($(current page.north east)+(0,-1.15cm)$);
    \fill[cAccent]
      (current page.north west)
      rectangle ($(current page.north west)+(0.28\paperwidth,-1.15cm)$);
    \node[white, anchor=west, font=\small\itshape] at
      ($(current page.north west)+(0.28\paperwidth+8pt,-0.58cm)$)
      {Go Ising \enspace\textbullet\enspace
       Experimentos~07+08 \enspace\textbullet\enspace
       Moyo y teor\'ia de invariantes};
  \end{tikzpicture}%
}

\fancyfoot[L]{\small\color{cMuted}
  Leonardo Jim\'enez Mart\'inez \enspace\textbullet\enspace UNAM \enspace\textbullet\enspace 2026}
\fancyfoot[R]{\small\color{cMuted} P\'agina \thepage}
\def\footrule{{\color{cRule}\vspace{-3pt}\hrule width\headwidth height 0.4pt\vspace{2pt}}}

% -- Listas ----------------------------------------------------------------------------
\usepackage{enumitem}
\setlist[itemize]{
  leftmargin = 1.5em, itemsep = 2pt, topsep = 4pt,
  label = {\color{cAccent}\small$\bullet$}
}
\setlist[enumerate]{leftmargin=1.8em, itemsep=2pt, topsep=4pt}

% -- PDF metadata ------------------------------------------------------------------------
\usepackage[
  colorlinks = true,
  linkcolor  = cAccent,
  urlcolor   = cAccent,
  pdfauthor  = {Leonardo Jim\'enez Mart\'inez},
  pdftitle   = {Predicci\'on de moyos con Hamiltonianos Go -- Informe completo (evidencia emp\'irica y teor\'ia de invariantes)},
  pdfsubject = {Go Ising -- Experimentos 07 y 08}
]{hyperref}

% -- Macros -------------------------------------------------------------------------------
\newcommand{\TH}[1]{\cellcolor{cTHbg}\color{white}\textbf{#1}}
\newcommand{\good}[1]{\color{cGood}\textbf{#1}}
\newcommand{\bad}[1]{\color{cBad}\textbf{#1}}

% =================================================================================
\begin{document}
\thispagestyle{fancy}

% -- Portada ----------------------------------------------------------------------------
\vspace*{4pt}
{\color{cNavy}\Huge\bfseries Predicci\'on de moyos}\\[3pt]
{\color{cNavy}\Large\bfseries con Hamiltonianos cl\'asicos de Ising}

\vspace{6pt}
{\color{cMuted}\large\itshape
  Informe completo --- evidencia emp\'irica (Experimento~07) y
  fundamento te\'orico (Experimento~08)}

\vspace{4pt}
{\color{cMuted}\small\itshape Generado autom\'aticamente \textbullet\ 2026-07-19}\\[1pt]
{\color{cMuted}\small\itshape
  20 partidas profesionales reales \textbullet\ motor de an\'alisis KataGo local
  (kata1-b15c192) \textbullet\ verificaci\'on simb\'olica con \texttt{sympy}}

\vspace{8pt}
{\color{cAccent}\hrule height 2pt}
\vspace{10pt}

% =================================================================================
\section{C\'omo se relacionan las dos partes de este informe}
\label{sec:relacion}

Este documento une dos experimentos de naturaleza distinta que, sin
embargo, son la misma investigaci\'on vista desde dos \'angulos:

\begin{itemize}
  \item \textbf{Parte~I (Experimento~07) es emp\'irica.} Usa
  Hamiltonianos cl\'asicos de Ising y un campo de relajaci\'on de
  Boltzmann para \textbf{predecir} el moyo/territorio real que produce
  KataGo al analizar partidas profesionales. Corre KataGo, mide
  $\Delta R^2$, compara candidatos --- es estad\'istica aplicada sobre
  datos reales.

  \item \textbf{Parte~II (Experimento~08) es te\'orica.} No predice
  nada nuevo ni corre KataGo. Toma \textbf{un solo hallazgo puntual}
  que surgi\'o dentro de la Parte~I (Secci\'on~\ref{sec:degeneracion}:
  ``el campo de relajaci\'on solo depende de 4 de los 7 coeficientes de
  \texttt{cubic\_mixed}, nunca de los 7 por separado'') y explica
  \textbf{por qu\'e} ocurre eso, de forma sistem\'atica, usando teor\'ia
  de grupos (el grupo de Klein). Es la justificaci\'on formal de un
  mecanismo que la Parte~I ya estaba usando sin demostrar.
\end{itemize}

Dicho de otro modo: la Parte~I responde ``¿qu\'e tan bien predice este
Hamiltoniano el moyo real?''; la Parte~II responde ``¿por qu\'e el
campo de relajaci\'on ignora ciertos coeficientes del Hamiltoniano, sin
importar cu\'al Hamiltoniano sea?''. La segunda pregunta naci\'o
\textbf{dentro} del primer experimento, pero su respuesta es puramente
algebraica --- no depende de correr KataGo ni de datos --- por lo que
se independiz\'o en su propio experimento con su propia carpeta
(\texttt{experiments/08\_teoria\_invariantes/}). Este informe las
presenta juntas, en el orden en que l\'ogicamente se necesitan: primero
la evidencia que motiva la pregunta (Parte~I), despu\'es la respuesta
formal (Parte~II).

% =================================================================================
\part{Predicci\'on emp\'irica de moyo}
\label{part:07}

% =================================================================================
\section{Objetivo}

Determinar si un Hamiltoniano cl\'asico de interacci\'on por pares
$H(s_i,s_j)$, definido sobre el modelo de esp\'in de Go
$s\in\{-1,0,+1\}$, puede predecir el \textbf{territorio real} (\emph{moyo})
que produce un motor de IA de referencia (KataGo) al analizar partidas
profesionales, mejor de lo que predice la sola geometr\'ia del tablero
(distancia a piedras, distancia al borde).

La m\'etrica de \'exito es $\Delta R^2$: la ganancia de $R^2$ al agregar
el campo promedio del Hamiltoniano sobre una regi\'on candidata,
por encima de un modelo base que solo usa geometr\'ia, con significancia
verificada por prueba $F$.

% =================================================================================
\section{Flujo de trabajo}

El pipeline separa tres etapas; \textbf{KataGo solo interviene en la
segunda}, y se corre una sola vez por posici\'on:

\begin{enumerate}
  \item \textbf{Partida $\to$ tablero.} Se parsea el \texttt{.sgf}, se
  muestrean $\sim$6 momentos por partida y se reconstruye el tablero
  $19\times19$ en cada uno, antes de invocar a KataGo.

  \item \textbf{An\'alisis KataGo.} \texttt{katago.exe} se lanza como
  proceso persistente en modo \texttt{analysis} (protocolo JSON por
  \texttt{stdin}/\texttt{stdout}, backend OpenCL sobre GPU). Por cada
  posici\'on se env\'ia una consulta con \texttt{maxVisits} fijo y se
  recibe \texttt{ownership} (estimaci\'on de control por punto) y
  \texttt{rootInfo.scoreLead}. Es \textbf{sin estado}: cada proceso es
  independiente, no existe base de datos persistente, y el motor mezcla
  aleatoriedad interna entre corridas separadas.

  \item \textbf{Post-proceso (sin KataGo).} Con el \texttt{ownership}
  crudo se detectan regiones de moyo/territorio (\emph{flood-fill}
  por bandas de control), se calculan features geom\'etricas, y por
  separado se eval\'ua el campo de relajaci\'on de Boltzmann del
  Hamiltoniano candidato --- este \'ultimo paso no requiere a KataGo,
  es puro c\'alculo del modelo Ising cl\'asico.
\end{enumerate}

\begin{tcolorbox}[quotebox]
  \textbf{Decisi\'on de dise\~no confirmada}: correr KataGo \emph{una
  sola vez} por posici\'on y cachear tablero + regiones + geometr\'ia
  en disco, reutilizando ese cache para evaluar cientos de
  Hamiltonianos distintos (solo cambia el c\'alculo barato del campo).
  Es el patr\'on que hace viable la b\'usqueda de coeficientes descrita
  en la Secci\'on~\ref{sec:optimizacion}.
\end{tcolorbox}

% =================================================================================
\section{Qu\'e es exactamente $H(x,y)$, y c\'omo se convierte en una predicci\'on}
\label{sec:que-es-h}

Esta secci\'on responde, de forma autocontenida, dos preguntas que no
son obvias a partir de los resultados solos: \textbf{qu\'e es} el
Hamiltoniano que se est\'a probando, y \textbf{por qu\'e} hace falta un
proceso adicional (relajaci\'on de Boltzmann) para convertirlo en una
predicci\'on, en vez de simplemente "calcular su energ\'ia".

\subsection{$H(x,y)$ es un Hamiltoniano de Ising cl\'asico, generalizado}

En el modelo de Ising cl\'asico de libro de texto, la energ\'ia total
del sistema es una suma de t\'erminos de interacci\'on entre pares de
espines vecinos:
\[
H_{\text{total}} = \sum_{\langle i,j\rangle \text{ vecinos}} H(s_i,s_j),
\qquad H(s_i,s_j) = -J\,s_i s_j \quad \text{(caso cl\'asico est\'andar)}.
\]
Ese t\'ermino local, $-Js_is_j$, es \'el mismo un polinomio (bilineal,
grado 2). Este proyecto generaliza \'unicamente esa pieza: en vez de
fijar la forma del acoplamiento a $-Js_is_j$, se permite que $H(x,y)$
sea cualquier polinomio cont\'inuo de dos variables --- por ejemplo,
la plantilla \texttt{cubic\_mixed},
$H(x,y)=a_1x+a_2y+b_{11}x^2+b_{12}xy+b_{22}y^2+c_{112}x^2y+c_{122}xy^2$,
de la cual $H_{M1}$, Alvarado, $H_{OPT\_A}$ y $H_{OPT\_B}$ son casos
particulares con coeficientes espec\'ificos (Secci\'on~\ref{sec:los7}
los detalla todos). Sigue siendo, formalmente, un Hamiltoniano de Ising
cl\'asico --- una suma de energ\'ias de interacci\'on par a par sobre
el tablero --- solo que con una funci\'on de acoplamiento m\'as rica
que el simple producto $s_is_j$. El espacio de esp\'ines es
$s\in\{-1,0,+1\}$ (blanco, vac\'io, negro); el modelo es enteramente
\textbf{cl\'asico} en todo el proyecto, nunca cu\'antico.

\subsection{Por qu\'e "solo energ\'ia" no alcanza}
\label{sec:por-que-boltzmann}

$H(x,y)$ punt\'ua la interacci\'on de un \textbf{par} de espines --- no
est\'a definida para un solo punto. Un punto vac\'io del tablero
todav\'ia no tiene un valor de esp\'in asignado: eso es justo lo que se
quiere averiguar. Si se intentara usar la energ\'ia directamente, sin
ning\'un proceso adicional, aparecen tres huecos que la f\'ormula sola
no resuelve:

\begin{enumerate}
  \item \textbf{No hay un n\'umero \'unico por punto.} Para un punto
  vac\'io $p$, se pueden probar distintos valores hipot\'eticos de
  esp\'in $s\in[-1,1]$ y calcular la energ\'ia de interacci\'on con
  cada vecino para cada uno --- eso da una \textbf{curva completa} de
  energ\'ias (una por cada $s$ candidato), no un solo valor resumido
  que se pueda comparar contra el \texttt{pct\_black} real (que s\'i es
  un solo n\'umero por regi\'on).
  \item \textbf{No hay propagaci\'on.} $H(s,q)$ solo conecta un punto
  con sus vecinos dentro del radio de interacci\'on, en un solo paso.
  El territorio real de una partida depende de piedras que pueden estar
  m\'as lejos que ese radio --- hace falta que la influencia de una
  piedra se propague a trav\'es de varios puntos vac\'ios intermedios.
  \item \textbf{No hay una regla de agregaci\'on con sentido f\'isico.}
  Dado el conjunto de energ\'ias candidatas por punto, ¿cu\'al valor de
  esp\'in "gana"? Tomar el m\'inimo ser\'ia una elecci\'on arbitraria y
  discreta (con empates mal definidos).
\end{enumerate}

\subsection{El proceso de relajaci\'on de Boltzmann}
\label{sec:mecanismo-relajacion}

La predicci\'on f\'isicamente correcta, en mec\'anica estad\'istica
cl\'asica, es el \textbf{valor esperado} del esp\'in de cada punto bajo
la distribuci\'on de Boltzmann que el Hamiltoniano define:
\[
\langle s_p \rangle = \frac{\sum_{\text{configuraciones}} s_p\cdot
e^{-H_{\text{total}}/T}}{\sum_{\text{configuraciones}} e^{-H_{\text{total}}/T}}.
\]
Calcular eso exactamente sobre un tablero $19\times19$ es intratable
(demasiadas configuraciones posibles). \texttt{relaxation\_field}
(\texttt{features.py:81-141}) aproxima esa misma cantidad con un
m\'etodo de relajaci\'on iterativo, est\'andar en modelos de Ising
cuando la soluci\'on exacta no es viable: para cada punto vac\'io $p$,
se prueban 41 valores de esp\'in candidatos en $[-1,1]$, se calcula la
energ\'ia local
\[
E_{\text{local}}(s) = \sum_{q\ \text{vecino de}\ p} \text{coef}(d_q)\cdot
\bigl[H(s,F[q]) + H(F[q],s)\bigr],
\]
(la suma sim\'etrica $H(s,q)+H(q,s)$ es la que motiva la reducci\'on
del grupo de Klein, Parte~\ref{part:08}) y se promedia con peso de
Boltzmann:
\[
F[p] = \frac{\sum_s s\cdot e^{-E_{\text{local}}(s)/T}}{\sum_s e^{-E_{\text{local}}(s)/T}}.
\]
Esto se repite \texttt{n\_sweeps} veces (barridos): cada barrido deja
que la influencia avance un paso m\'as, propag\'andose m\'as all\'a del
radio de interacci\'on de un solo paso. Las piedras ya puestas quedan
\textbf{fijas} durante todo el proceso (condiciones de frontera, no se
relajan). Este procedimiento no es un reemplazo del Hamiltoniano ---
es c\'omo se extrae de \'el una predicci\'on observable, dado que la
soluci\'on exacta de mec\'anica estad\'istica es computacionalmente
inviable.

\subsection{Qu\'e es \texttt{H\_field\_mean}}

Es la cantidad que aparece en todas las tablas de $\Delta R^2$ del
informe. \textbf{No} es $H(x,y)$ evaluado directamente en un punto ---
es el promedio del campo $F$ ya relajado (tras los \texttt{n\_sweeps}
barridos completos) sobre todos los puntos de una regi\'on candidata
(un moyo, un territorio, una esquina de joseki, etc.):
\[
\texttt{H\_field\_mean} = \frac{1}{|\text{regi\'on}|}\sum_{p\in\text{regi\'on}} F[p].
\]
Es esta cantidad --- no $H$ ni la energ\'ia cruda --- la que se agrega
como predictor en la regresi\'on que produce cada $\Delta R^2$
reportado en este documento.

% =================================================================================
\section{Lo que ya se descart\'o}

\subsection{Criterios topol\'ogicos (Frente 1 de Pareto)}

El experimento~06 hab\'ia ordenado los 44 candidatos \texttt{cubic\_mixed}
por criterios topol\'ogicos (persistencia $H_1$, separaci\'on de puntos
cr\'iticos, rango energ\'etico), definiendo un ``Frente~1'' de candidatos
te\'oricamente superiores. Contrastado contra $\Delta R^2$ real:

\medskip
\begin{center}
\rowcolors{2}{white}{cTRalt}
\begin{tabular}{lrr}
  \TH{Grupo (Hasse)} & \TH{$\overline{\Delta R^2}$ (r=1)} & \TH{$\overline{\Delta R^2}$ (r=9)} \\[1pt]
  Frente 1 (``mejor'' te\'orico) & 0.130 & \bad{0.055} \\
  Medios                          & 0.204 & 0.170 \\
  Tard\'ios                       & 0.216 & 0.152 \\
\end{tabular}
\end{center}
\medskip

El Frente~1 --- el grupo que la topolog\'ia predec\'ia como superior ---
resulta ser el de \textbf{peor} desempe\~no real, y el que m\'as se
degrada al ampliar el radio de interacci\'on. Los criterios topol\'ogicos
no predicen utilidad real para este problema.

\subsection{Muestreo aleatorio, en cualquier tama\~no de familia}

Se probaron dos familias de Hamiltonianos generados al azar
(\texttt{seed} fija):

\medskip
\begin{center}
\rowcolors{2}{white}{cTRalt}
\begin{tabularx}{0.96\textwidth}{L{3.4cm}C{2cm}C{2.3cm}C{2.3cm}X}
  \TH{Familia} & \TH{Candidatos} & \TH{$\overline{\Delta R^2}$ r=1} & \TH{$\overline{\Delta R^2}$ r=9} & \TH{Tendencia con radio} \\[1pt]
  \texttt{cubic\_mixed} (7 par\'am.) & 44 & 0.130--0.216 & 0.055--0.170 & \bad{degrada} \\
  \texttt{sparse\_cubic} (4 par\'am., misma forma que $H_{M1}$) & 25 & 0.102 & 0.056 & \bad{degrada} \\
\end{tabularx}
\end{center}
\medskip

Reducir la familia a los mismos 4 mon\'omios de $H_{M1}$ y muestrear
sus coeficientes al azar \textbf{no ayuda}: el promedio (0.056--0.102)
es incluso peor que el de la familia completa de 7 par\'ametros. Lo que
importa no es la \emph{forma} de los mon\'omios sino el \emph{valor}
espec\'ifico de sus coeficientes.

\subsection{Inestabilidad de los mejores candidatos aleatorios}

El mejor candidato aleatorio individual, $H_{0106}$, ilustra el
problema con m\'as claridad que cualquier promedio:

\medskip
\begin{center}
\rowcolors{2}{white}{cTRalt}
\begin{tabular}{lrrr}
  \TH{Hamiltoniano} & \TH{$\Delta R^2$ (r=1)} & \TH{$\Delta R^2$ (r=9)} & \TH{Cambio} \\[1pt]
  $H_{0106}$ (aleatorio, mejor en r=1) & \good{0.408} & \bad{0.005} & $-0.403$ \\
  $H_{0115}$ (aleatorio, mejor en r=9) & 0.385 & \good{0.437} & $+0.052$ \\
\end{tabular}
\end{center}
\medskip

$H_{0106}$ es el mejor candidato aleatorio a radio~1 y se
\textbf{colapsa} a radio~9 --- una ca\'ida de m\'as del 98\,\%. Ning\'un
candidato aleatorio es consistentemente bueno en ambos reg\'imenes.

% =================================================================================
\section{Lo que s\'i funciona: los Hamiltonianos derivados a mano}

$H_{M1}$ (Mercado \& Jim\'enez) y el modelo de Alvarado (Atomic-Go) no
fueron generados al azar --- sus coeficientes vienen de una construcci\'on
te\'orica previa, ajena a este pipeline. Son, hasta ahora, los
\textbf{\'unicos} Hamiltonianos probados que \textbf{mejoran} al ampliar
el radio de interacci\'on de Manhattan:

\medskip
\begin{center}
\rowcolors{2}{white}{cTRalt}
\begin{tabular}{lrrr}
  \TH{Hamiltoniano} & \TH{$\Delta R^2$ (r=1)} & \TH{$\Delta R^2$ (r=9)} & \TH{Cambio} \\[1pt]
  $H_{M1}$ (\texttt{sparse\_cubic}) & 0.242 & \good{0.309} & \good{$+0.067$} \\
  Alvarado (\texttt{quadratic})     & 0.169 & \good{0.255} & \good{$+0.086$} \\
\end{tabular}
\end{center}
\medskip

\begin{tcolorbox}[quotebox]
  Este es el hallazgo emp\'irico que motiva el pivote de la
  Secci\'on~\ref{sec:optimizacion}: la mejora sostenida con el radio no
  parece ser un accidente del muestreo, sino una propiedad de
  coeficientes \emph{deliberadamente} elegidos. Ninguno de los 69
  candidatos generados al azar (44 + 25) replica ese comportamiento.
\end{tcolorbox}

% =================================================================================
\section{Pivote: optimizaci\'on directa de coeficientes}
\label{sec:optimizacion}

Dado que ni el criterio topol\'ogico ni el muestreo aleatorio producen
candidatos consistentemente buenos, se reemplaz\'o la b\'usqueda al
azar por \textbf{optimizaci\'on directa} contra la se\~nal real
($\Delta R^2$), usando \texttt{scipy.optimize.differential\_evolution}.

\subsection{Infraestructura de cach\'e}

\begin{tcolorbox}[codebox]
\texttt{cache\_positions.py} \\
\ \ \ $\to$ corre KataGo 1 vez por posici\'on \\
\ \ \ $\to$ guarda tablero + moyos detectados + features geom\'etricas \\[4pt]
\texttt{optimize\_coefficients.py} \\
\ \ \ $\to$ reutiliza el cache cientos de veces \\
\ \ \ $\to$ solo recalcula el campo (barato) para cada vector de coeficientes
\end{tcolorbox}

La funci\'on objetivo se verific\'o antes de optimizar: evaluada en los
coeficientes exactos de $H_{M1}$ sobre un cache de 6 partidas, reproduce
$\Delta R^2 \approx 0.315$--$0.329$ (seg\'un n\'umero de barridos de
relajaci\'on), consistente con el $0.309$ medido sobre las 20 partidas
completas.

\subsection{Fase A --- prueba de concepto}

\begin{center}
\rowcolors{2}{white}{cTRalt}
\begin{tabular}{ll}
  \TH{Par\'ametro} & \TH{Valor} \\[1pt]
  Plantilla & \texttt{sparse\_cubic} (4 coeficientes, forma de $H_{M1}$) \\
  Radio de Manhattan & 9 \\
  Partidas en cach\'e & 6 \\
  Arranques sembrados & $H_{M1}$, y proyecciones de $H_{0106}$/$H_{0115}$ \\
  Optimizador & \texttt{differential\_evolution} (\texttt{workers=1}) \\
\end{tabular}
\end{center}

\textbf{Resultado} (547 evaluaciones, $\sim$75 min, sobre el cach\'e
de 6 partidas):

\medskip
\begin{center}
\rowcolors{2}{white}{cTRalt}
\begin{tabular}{ll}
  \TH{Coeficiente} & \TH{Valor \'optimo} \\[1pt]
  $a_1$    & $-0.8404$ \\
  $a_2$    & $-0.1608$ \\
  $c_{112}$ & $\phantom{-}0.9828$ \\
  $c_{122}$ & $-0.1221$ \\
  $\Delta R^2$ (cach\'e de 6 partidas) & \good{0.3523} \\
\end{tabular}
\end{center}
\medskip

Superior al $\Delta R^2=0.3152$ de $H_{M1}$ real medido sobre el mismo
cach\'e y mismas condiciones ($n\_sweeps=4$). El optimizador no
convergi\'o formalmente (\texttt{success: false}, la funci\'on
objetivo se aplan\'o hacia el final --- ver Secci\'on~\ref{sec:degeneracion}
sobre por qu\'e el paisaje tiene direcciones exactamente planas), pero
el resultado final s\'i mejora a la referencia.

\subsubsection{Verificaci\'on cruzada con la identidad de degeneraci\'on}

Sumando los coeficientes \'optimos: $\Sigma a = a_1+a_2 = -1.0013$,
$\Sigma c = c_{112}+c_{122} = 0.8607$. Seg\'un la
Secci\'on~\ref{sec:degeneracion}, cualquier otro reparto de esas
mismas sumas deber\'ia dar id\'entico $\Delta R^2$. Se prob\'o un
reparto alterno, $(a_1,a_2,c_{112},c_{122}) = (2.499,-3.501,-0.270,1.130)$,
misma $\Sigma a$ y $\Sigma c$ pero valores individuales muy distintos:

\begin{tcolorbox}[quotebox]
  \'Optimo original: $\Delta R^2=0.3523$. Reparto alterno:
  $\Delta R^2=0.3523$. \textbf{Id\'entico a 4 cifras decimales} ---
  confirma emp\'iricamente, sobre un punto real hallado por el
  optimizador (no solo sobre coeficientes unitarios sint\'eticos),
  que la identidad algebraica de la Secci\'on~\ref{sec:degeneracion}
  es exacta. La demostraci\'on formal de \textbf{por qu\'e} esta
  identidad tiene que cumplirse --- no solo que se cumple --- es el
  tema completo de la Parte~\ref{part:08}.
\end{tcolorbox}

\subsubsection{Validaci\'on sobre las 20 partidas completas}

Aplicando el est\'andar de rigor de la Secci\'on~\ref{sec:barrido9}
(nunca aceptar un candidato sin probarlo fuera del subset donde se
encontr\'o), se evalu\'o el \'optimo de Fase~A sobre el cach\'e
completo de 20 partidas:

\medskip
\begin{center}
\rowcolors{2}{white}{cTRalt}
\begin{tabular}{lrr}
  \TH{Hamiltoniano} & \TH{$\Delta R^2$ ($n{=}4$ barridos)} & \TH{$\Delta R^2$ ($n{=}8$ barridos)} \\[1pt]
  \'Optimo de Fase~A & \good{0.3112} & \good{0.3165} \\
  $H_{M1}$ real       & 0.2858 & 0.2981 \\
\end{tabular}
\end{center}
\medskip

\begin{tcolorbox}[quotebox]
  \textbf{A diferencia de la clase espejo de la
  Secci\'on~\ref{sec:barrido9}, esta mejora s\'i se sostiene} al pasar
  de 6 a 20 partidas, en ambos ajustes de barridos. Es la primera vez
  en este experimento que un Hamiltoniano supera consistentemente a
  $H_{M1}$ sobre el cach\'e completo, y sali\'o de b\'usqueda
  deliberada, no de muestreo al azar ni de una coincidencia de subset
  peque\~no.
\end{tcolorbox}

\subsubsection{Tambi\'en mejora con el radio, m\'as que $H_{M1}$}

La Secci\'on~4 identific\'o como propiedad distintiva de $H_{M1}$ y
Alvarado que su $\Delta R^2$ \textbf{mejora} al ampliar el radio de
Manhattan --- ning\'un candidato aleatorio de los 69 probados hac\'ia
eso. El \'optimo de Fase~A, evaluado en ambos radios sobre las 20
partidas completas ($n\_sweeps=8$, id\'entico a la Secci\'on~4):

\medskip
\begin{center}
\rowcolors{2}{white}{cTRalt}
\begin{tabular}{lrrr}
  \TH{Hamiltoniano} & \TH{$\Delta R^2$ (r=1)} & \TH{$\Delta R^2$ (r=9)} & \TH{Cambio} \\[1pt]
  \'Optimo de Fase~A & 0.235 & \good{0.3165} & \good{$+0.082$} \\
  $H_{M1}$ real       & 0.242 & 0.2981          & $+0.056$ \\
\end{tabular}
\end{center}
\medskip

No solo supera a $H_{M1}$ en ambos radios: tambi\'en mejora \emph{m\'as}
al ampliar el radio ($+0.082$ contra $+0.056$). Es el tercer
Hamiltoniano --- y el primero encontrado por b\'usqueda deliberada en
vez de derivaci\'on te\'orica externa --- que exhibe esta propiedad.

\subsection{Degeneraci\'on algebraica de \texttt{sparse\_cubic}}
\label{sec:degeneracion}

Mientras se esperaba el resultado de Fase A se audit\'o la
implementaci\'on de \texttt{relaxation\_field} (\texttt{features.py:136})
y se encontr\'o una identidad exacta que reduce el espacio de
b\'usqueda efectivo. El campo nunca eval\'ua $H(s,q)$ sola: cada
actualizaci\'on usa siempre la suma sim\'etrica $H(s,q)+H(q,s)$, donde
$s$ es el punto que se actualiza y $q$ cada vecino en el kernel:

\begin{tcolorbox}[codebox]
\texttt{e\_matrix = CF * (h(S, NV) + h(NV, S))}
\end{tcolorbox}

Para la plantilla \texttt{sparse\_cubic},
$H(x,y)=a_1x+a_2y+c_{112}x^2y+c_{122}xy^2$, esa suma se simplifica
algebraicamente (v\'alido para $s,q\in\mathbb{R}$, no solo en
$\{-1,0,1\}$):

\begin{tcolorbox}[codebox]
$H(s,q)+H(q,s) \;=\; (a_1+a_2)(s+q) \;+\; (c_{112}+c_{122})\,(s^2q+sq^2)$
\end{tcolorbox}

\begin{tcolorbox}[quotebox]
  \textbf{Consecuencia}: el campo de relajaci\'on --- y por lo tanto
  todo el $\Delta R^2$ que se est\'a optimizando --- depende
  \'unicamente de las dos sumas $\Sigma a = a_1+a_2$ y
  $\Sigma c = c_{112}+c_{122}$, nunca de $a_1,a_2,c_{112},c_{122}$
  por separado. El espacio de b\'usqueda de 4 par\'ametros contin\'uos
  de Fase~A tiene direcciones exactamente planas: infinitos pares
  $(a_1,a_2)$ con la misma suma --- e infinitos pares
  $(c_{112},c_{122})$ con la misma suma --- producen exactamente el
  mismo $\Delta R^2$. \textbf{Por qu\'e tiene que ser exactamente as\'i}
  (y no una coincidencia de esta plantilla en particular) se deriva
  formalmente en la Parte~\ref{part:08}.
\end{tcolorbox}

\subsubsection{Clasificaci\'on con coeficientes unitarios}

Restringiendo cada uno de los 4 par\'ametros a $\pm1$ hay
$2^4=16$ combinaciones de signos posibles. Bajo la identidad anterior,
$\Sigma a\in\{-2,0,+2\}$ y $\Sigma c\in\{-2,0,+2\}$, cada uno con
multiplicidad $1,2,1$ seg\'un el signo --- lo que reduce las 16
combinaciones a solo \textbf{9 clases efectivas distintas}:

\medskip
\begin{center}
\rowcolors{2}{white}{cTRalt}
\begin{tabular}{ccc}
  \TH{$\Sigma a=a_1+a_2$} & \TH{$\Sigma c=c_{112}+c_{122}$} & \TH{Miembros (de 16) en esta clase} \\[1pt]
  $-2$ & $-2$ & 1 \\
  $-2$ & $0$  & 2 \\
  $-2$ & $+2$ & 1 \\
  $0$  & $-2$ & 2 \\
  $0$  & $0$  & \bad{4 (m\'axima redundancia)} \\
  $0$  & $+2$ & 2 \\
  $+2$ & $-2$ & \good{1 (clase de $H_{M1}$)} \\
  $+2$ & $0$  & 2 \\
  $+2$ & $+2$ & 1 \\
\end{tabular}
\end{center}
\medskip

$H_{M1}$ tiene signos $(+,+,-,-)$, que cae en la clase
$(\Sigma a{=}2,\ \Sigma c{=}{-}2)$ --- una clase \emph{singleton}: no
hay otro patr\'on de signos unitarios que produzca el mismo campo.

\begin{tcolorbox}[pendingbox]
  \textbf{Advertencia de alcance}: esta reducci\'on es exclusiva de
  \texttt{relaxation\_field} (el pipeline de predicci\'on de moyo). No
  aplica al an\'alisis topol\'ogico de $\Gamma(H)$ del experimento~06
  (nodos $A_1$, puntos cr\'iticos, g\'enero de la fibra), que s\'i
  distingue $a_1$ de $a_2$ y $c_{112}$ de $c_{122}$ individualmente ---
  ese an\'alisis eval\'ua $H(x,y)$ directamente, sin simetrizar.
\end{tcolorbox}

\textbf{Implicaci\'on pr\'actica}: para Fase~B, reparametrizar la
b\'usqueda de \texttt{sparse\_cubic} como 2 par\'ametros efectivos
($\Sigma a$, $\Sigma c$) en vez de 4 evitar\'ia gastar presupuesto de
evaluaci\'on en direcciones planas del objetivo --- estrictamente m\'as
eficiente sin p\'erdida, \emph{solo} para el objetivo de $\Delta R^2$.

\subsection{Barrido ordenado de las 9 clases efectivas}
\label{sec:barrido9}

Antes de seguir esperando a la optimizaci\'on continua de Fase~A, se
aprovech\'o la reducci\'on de la Secci\'on~\ref{sec:degeneracion} para
evaluar de forma \textbf{exhaustiva y determinista} --- sin
aleatoriedad, sin optimizador --- las 9 clases efectivas con
coeficientes unitarios, contra el mismo cach\'e de 6 partidas de
Fase~A (radio~9, 4 barridos de relajaci\'on):

\medskip
\begin{center}
\rowcolors{2}{white}{cTRalt}
\begin{tabular}{rrcp{4cm}rr}
  \TH{$\Sigma a$} & \TH{$\Sigma c$} & \TH{miembros} & \TH{representante $(a_1,a_2,c_{112},c_{122})$} & \TH{$\Delta R^2$} & \TH{$p$-valor} \\[1pt]
  $-2$ & $+2$ & 1 & $(-1,-1,+1,+1)$ & \good{0.3263} & $1.1\times10^{-16}$ \\
  $+2$ & $-2$ & 1 & $(+1,+1,-1,-1)$ \small (signos de $H_{M1}$) & 0.3013 & $1.1\times10^{-16}$ \\
  $0$  & $+2$ & 2 & $(-1,+1,+1,+1)$ & 0.1019 & $2.1\times10^{-10}$ \\
  $+2$ & $+2$ & 1 & $(+1,+1,+1,+1)$ & 0.0671 & $3.6\times10^{-7}$ \\
  $0$  & $-2$ & 2 & $(-1,+1,-1,-1)$ & 0.0529 & $7.1\times10^{-6}$ \\
  $-2$ & $-2$ & 1 & $(-1,-1,-1,-1)$ & 0.0277 & $1.3\times10^{-3}$ \\
  $-2$ & $0$  & 2 & $(-1,-1,-1,+1)$ & 0.0044 & $0.201$ \\
  $+2$ & $0$  & 2 & $(+1,+1,-1,+1)$ & 0.0044 & $0.201$ \\
  $0$  & $0$  & 4 & $(-1,+1,-1,+1)$ & \bad{0.0000} & $1.000$ \\
\end{tabular}
\end{center}
\medskip

Dos confirmaciones y un hallazgo nuevo:

\begin{itemize}
  \item \textbf{Confirma la identidad algebraica}: la clase
        $(\Sigma a{=}0,\Sigma c{=}0)$ da $\Delta R^2=0.0000$ exacto
        --- el campo simetrizado es id\'enticamente cero para
        cualquiera de sus 4 miembros, como predice la
        Secci\'on~\ref{sec:degeneracion}.
  \item \textbf{Confirma a $H_{M1}$}: la clase con el mismo patr\'on de
        signos que $H_{M1}$ ($+,+,-,-$) da $\Delta R^2=0.3013$,
        consistente con el $0.309$ medido sobre las 20 partidas
        completas (Secci\'on 4).
  \item \textbf{Candidato provisional}: en el subset de 6 partidas, la
        clase \emph{espejo} de $H_{M1}$ --- todos los signos
        invertidos, $(-1,-1,+1,+1)$ --- rendia \emph{mejor} que la
        clase de $H_{M1}$: $0.3263$ contra $0.3013$.
\end{itemize}

\begin{tcolorbox}[quotebox]
  \textbf{Verificaci\'on sobre las 20 partidas completas} (115
  posiciones, mismo radio y barridos): la ventaja del espejo
  \textbf{no se sostiene}. Con el cach\'e completo, la clase de
  $H_{M1}$ vuelve a ganar:

  \medskip
  \begin{center}
  \rowcolors{2}{white}{cTRalt}
  \begin{tabular}{lrr}
    \TH{Clase} & \TH{$\Delta R^2$ (6 partidas)} & \TH{$\Delta R^2$ (20 partidas)} \\[1pt]
    $H_{M1}$ $(+1,+1,-1,-1)$ & 0.3013 & \good{0.3158} \\
    Espejo $(-1,-1,+1,+1)$   & \good{0.3263} & 0.3034 \\
  \end{tabular}
  \end{center}
  \medskip

  El orden se invierte al pasar de 6 a 20 partidas: lo que parec\'ia
  un hallazgo nuevo era ruido de muestra peque\~na, no una propiedad
  real del espejo. Es exactamente el mismo tipo de inestabilidad
  documentado para $H_{0106}$ en la Secci\'on~3.3 --- la raz\'on por
  la que este informe exige confirmar sobre el cach\'e completo antes
  de aceptar cualquier candidato, incluso uno que surja de una
  clasificaci\'on ordenada y no de b\'usqueda al azar.
\end{tcolorbox}

El resto de las 9 clases mantiene el mismo orden relativo entre 6 y 20
partidas (la clase nula $\Sigma a=\Sigma c=0$ sigue en $\Delta R^2
\approx 0$ exacto), as\'i que la inversi\'on es espec\'ifica al par
$H_{M1}$/espejo, no un efecto general del tama\~no de muestra sobre
todas las clases.

% =================================================================================
\section{Explotando toda la informaci\'on de KataGo: resultados}
\label{sec:rich}

Se integraron al pipeline los dos campos identificados como
aprovechables (\texttt{ownershipStdev}, \texttt{policy}), y se subi\'o
\texttt{maxVisits} de 250 a \textbf{600} para reducir el ruido interno
de KataGo. Se reconstruy\'o el cach\'e de 20 partidas
(\texttt{cache\_full20\_rich}: 113 posiciones, 864 moyos) con estos
campos y se repiti\'o el an\'alisis de $H_{M1}$ y $H_{OPT\_A}$ en
cuatro variantes.

\subsection{Ponderar por incertidumbre de KataGo}

Hip\'otesis: parte de la varianza no explicada ($\sim$37\,\% en la
Secci\'on~\ref{sec:optimizacion}) podr\'ia ser incertidumbre \emph{del
propio KataGo} en zonas t\'acticamente no resueltas, no una limitaci\'on
del Hamiltoniano. Se pondera la regresi\'on por
$1/\texttt{ownership\_stdev\_mean}^2$ (m\'inimos cuadrados ponderados):

\medskip
\begin{center}
\rowcolors{2}{white}{cTRalt}
\begin{tabular}{lrr}
  \TH{Hamiltoniano} & \TH{$\Delta R^2$ sin ponderar} & \TH{$\Delta R^2$ ponderado} \\[1pt]
  $H_{M1}$      & 0.2967 & \good{0.4041} \\
  $H_{OPT\_A}$  & 0.3073 & \good{0.3832} \\
\end{tabular}
\end{center}
\medskip

\begin{tcolorbox}[quotebox]
  Hip\'otesis confirmada: al bajarle peso a los moyos donde KataGo mismo
  est\'a inseguro, el $\Delta R^2$ sube fuertemente para ambos ($+36\,\%$
  y $+25\,\%$ respectivamente). Una parte real de lo "no explicado"
  s\'i era ruido de KataGo, no falta de poder predictivo del modelo.
  \textbf{Pero el orden se invierte}: ponderando, $H_{M1}$ supera a
  $H_{OPT\_A}$ (0.404 contra 0.383) --- lo opuesto al resultado sin
  ponderar. Tercera vez en el experimento que un ranking se invierte
  al cambiar la metodolog\'ia de medici\'on, no los datos.
\end{tcolorbox}

\subsection{¿Aporta \texttt{policy} algo independiente?}

Se prob\'o \texttt{policy\_mass} (masa de la red de pol\'itica dentro
de la regi\'on) como predictor, tanto agregado a la geometr\'ia base
como su aporte marginal solo:

\medskip
\begin{center}
\rowcolors{2}{white}{cTRalt}
\begin{tabular}{lr}
  \TH{Prueba} & \TH{Resultado} \\[1pt]
  $\Delta R^2$ de $H$ agregando \texttt{policy\_mass} a la base & sin cambio (0.297 / 0.308, igual que sin \texttt{policy\_mass}) \\
  $\Delta R^2$ de \texttt{policy\_mass} solo, sobre geometr\'ia & \bad{0.0007} ($p=0.36$, no significativo) \\
\end{tabular}
\end{center}
\medskip

\texttt{policy\_mass} \textbf{no aporta nada} en este pipeline ---
a diferencia de \texttt{ownershipStdev}, esta hip\'otesis no se
confirm\'o. El \texttt{ownership} ya captura lo que importa para
predecir moyo; la red de pol\'itica no agrega se\~nal independiente
\'util aqu\'i.

\subsection{Lo que faltaba de KataGo: \texttt{moveInfos}, \texttt{scoreStdev}, \texttt{winrate}, \texttt{scoreLead}}

Para completar la auditor\'ia de toda la informaci\'on disponible de
KataGo, se agreg\'o tambi\'en \texttt{moveInfos} (las jugadas candidatas
del motor en la ra\'iz) --- se calcul\'o \texttt{n\_top\_moves\_in\_region}:
cu\'antas de las 10 mejores jugadas candidatas caen dentro de cada
regi\'on (si el motor a\'un considera buenas jugadas ah\'i, el marco
no est\'a resuelto). Tambi\'en se probaron \texttt{rootInfo.scoreStdev},
\texttt{winrate} (volatilidad global de la posici\'on) y, en una
revisi\'on posterior que hab\'ia quedado pendiente,
\texttt{rootInfo.scoreLead} (diferencial de puntuaci\'on esperado ---
ya se guardaba en el cach\'e desde el principio, pero nunca se hab\'ia
probado como predictor):

\medskip
\begin{center}
\rowcolors{2}{white}{cTRalt}
\begin{tabular}{lrr}
  \TH{Se\~nal (aporte solo, sobre geometr\'ia)} & \TH{$\Delta R^2$} & \TH{$p$-valor} \\[1pt]
  \texttt{n\_top\_moves\_in\_region} & \bad{0.0002} & 0.61 \\
  \texttt{score\_stdev}              & \bad{0.0002} & 0.62 \\
  \texttt{winrate}                   & \bad{0.0013} & 0.20 \\
  \texttt{score\_lead}               & \bad{0.0015} & 0.16 \\
\end{tabular}
\end{center}
\medskip

\begin{tcolorbox}[quotebox]
  Ninguna de las cuatro aporta se\~nal significativa, ni cambia
  materialmente el $\Delta R^2$ de $H_{M1}$ (0.277 con o sin
  \texttt{score\_lead} como control) o $H_{OPT\_A}$ al agregarse como
  control. Con esto se completa la auditor\'ia de
  \texttt{rootInfo}/\texttt{ownership}/\texttt{ownershipStdev}/%
  \texttt{policy}/\texttt{moveInfos}: de los \textbf{7} campos
  evaluados en este experimento, \textbf{solo \texttt{ownershipStdev}}
  (ponderaci\'on) aport\'o algo real. El resto de la informaci\'on de
  KataGo, para esta pregunta espec\'ifica (predecir moyo con un campo
  Ising cl\'asico), es redundante con lo que ya captura
  \texttt{ownership}.
\end{tcolorbox}

\subsection{Por qu\'e el cat\'alogo "b\'asico" (r1/r9) nunca se recalcul\'o con estas se\~nales}
\label{sec:por-que-no-basico}

El visor interactivo distingue dos corridas: el cat\'alogo
\textbf{"b\'asico"} (los 44 candidatos \texttt{cubic\_mixed} originales
+ $H_{M1}$ + Alvarado + $H_{OPT\_A}$, con el selector de radio 1/9) y
las categor\'ias \textbf{"completas"} (moyo, territorio, fuseki,
joseki --- una sola corrida unificada de \texttt{cache\_positions.py}
con \texttt{maxVisits=600}, \texttt{ownershipStdev} y \texttt{policy}
activados). El cat\'alogo b\'asico es \textbf{anterior} a que estas
banderas existieran en el pipeline, as\'i que sus filas simplemente no
tienen \texttt{scoreStdev}, \texttt{winrate}, \texttt{policy(pass)},
\texttt{ownStdev}, \texttt{policyMass} ni \texttt{nTopMoves} --- no es
un error de carga de datos, es que esas columnas nunca se calcularon
para esa corrida.

Recalcular el cat\'alogo b\'asico con estas se\~nales requerir\'ia
volver a correr KataGo sobre las 20 partidas (con las banderas nuevas)
y recalcular \texttt{relaxation\_field} para los 47 Hamiltonianos en
\textbf{ambos} radios (1 y 9) --- del orden de 80--90 minutos de
c\'omputo, el doble de cualquier corrida individual de este informe.

\begin{tcolorbox}[quotebox]
  \textbf{No se hizo, y no hac\'ia falta}: la Secci\'on anterior ya
  responde la pregunta que esa recomputaci\'on estar\'ia respondiendo
  ---¿estas se\~nales aportan algo?--- de forma rigurosa (con bootstrap
  por partida) sobre un conjunto de Hamiltonianos distinto pero
  conceptualmente equivalente (\texttt{cache\_full20\_cats}/%
  \texttt{cache\_early20}). La respuesta ya es un \textbf{no}
  consistente en las 4 se\~nales probadas. Gastar 80--90 minutos de
  c\'omputo para reproducir la misma conclusi\'on nula sobre otro
  subconjunto de candidatos no cambiar\'ia nada --- ser\'ia repetir una
  pregunta ya cerrada, no responder una nueva. Por eso el visor,
  en vez de mostrar esas columnas vac\'ias (puro "—") para la
  categor\'ia b\'asica, simplemente las oculta ah\'i y las muestra
  \'unicamente donde s\'i existen.
\end{tcolorbox}

\subsection{La correcci\'on m\'as importante: bootstrap por partida}
\label{sec:bootstrap}

Como se discuti\'o antes de esta secci\'on, los p-valores reportados en
todo el experimento ($p\approx10^{-16}$) tratan cada \textbf{fila}
(moyo) como independiente, cuando en realidad las filas de una misma
partida est\'an correlacionadas --- pseudo-replicaci\'on cl\'asica. Se
remuestre\'o por \textbf{partida completa} (bootstrap, 500 r\'eplicas,
$n=20$ partidas) para obtener un intervalo de confianza honesto de
$\Delta R^2$:

\medskip
\begin{center}
\rowcolors{2}{white}{cTRalt}
\begin{tabular}{lrr}
  \TH{Hamiltoniano} & \TH{$\Delta R^2$ media (bootstrap)} & \TH{IC 95\,\%} \\[1pt]
  $H_{M1}$      & 0.294 & $[0.223,\ 0.358]$ \\
  $H_{OPT\_A}$  & 0.305 & $[0.227,\ 0.380]$ \\
\end{tabular}
\end{center}
\medskip

\begin{tcolorbox}[quotebox]
  \textbf{Los intervalos se traslapan casi por completo.} Con solo 20
  partidas, \textbf{no hay evidencia estad\'isticamente confiable de
  que $H_{OPT\_A}$ sea mejor que $H_{M1}$} --- la diferencia puntual
  (0.307 vs.\ 0.298) est\'a bien dentro del ruido de muestreo a este
  tama\~no. Esto no invalida el hallazgo de Fase A, pero s\'i cambia
  c\'omo debe presentarse: es un candidato \textbf{prometedor y
  consistente con ser al menos tan bueno como $H_{M1}$}, no un
  reemplazo confirmado.
\end{tcolorbox}

\textbf{Implicaci\'on directa}: la pregunta de "¿cu\'antas partidas
hacen falta?" (Secci\'on~\ref{sec:optimizacion}) ten\'ia una respuesta
concreta esperando --- con 20 partidas el ancho del IC~95\,\% es de
$\sim$0.13--0.15 en $\Delta R^2$, demasiado ancho para distinguir entre
los mejores candidatos actuales. Hacen falta m\'as partidas (o un
conjunto de validaci\'on separado) antes de declarar un ganador.

\subsection{Estabilidad de las familias aleatorias ante un segundo seed}

Las medias de \texttt{cubic\_mixed} (44 candidatos) y
\texttt{sparse\_cubic} (25 candidatos) reportadas en este experimento
se generaron con un \'unico seed cada una (42 y 7) --- una brecha de
dise\~no experimental se\~nalada expl\'icitamente: sin una segunda
r\'eplica, no se pod\'ia distinguir "as\'i se comporta la familia" de
"as\'i sali\'o este seed en particular". Se gener\'o una segunda
tanda (seeds 123 y 99) y se reevalu\'o sobre el mismo cach\'e
(\texttt{cache\_full20\_cats}):

\medskip
\begin{center}
\rowcolors{2}{white}{cTRalt}
\begin{tabular}{lrr}
  \TH{Familia (radio)} & \TH{Seed original} & \TH{Seed nuevo} \\[1pt]
  \texttt{cubic\_mixed} (r=1)  & $0.135\ [0.115,\,0.156]$ & $0.131\ [0.102,\,0.160]$ \\
  \texttt{cubic\_mixed} (r=9)  & $0.103\ [0.072,\,0.134]$ & $0.099\ [0.064,\,0.134]$ \\
  \texttt{sparse\_cubic} (r=1) & $0.086\ [0.055,\,0.118]$ & $0.092\ [0.070,\,0.115]$ \\
  \texttt{sparse\_cubic} (r=9) & $0.050\ [0.018,\,0.081]$ & $0.044\ [0.019,\,0.069]$ \\
\end{tabular}
\end{center}
\medskip

\begin{tcolorbox}[quotebox]
  Los intervalos se traslapan casi por completo en los 4 casos ---
  \textbf{la media de cada familia es estable ante un segundo seed
  independiente}. Las conclusiones cualitativas del experimento
  (aleatorio rinde muy por debajo de $H_{M1}$/$H_{OPT\_A}$; ambas
  familias se degradan con el radio) no son producto de la suerte de
  un seed particular --- se replican. Esta brecha de dise\~no queda
  cerrada.
\end{tcolorbox}

% =================================================================================
\section{M\'as all\'a del moyo: otras categor\'ias de KataGo}
\label{sec:categorias}

El moyo no es la \'unica categor\'ia que \texttt{detect\_moyos()} puede
dar. La funci\'on ya clasifica cada regi\'on vac\'ia en 3 bandas de
\texttt{ownership} --- \texttt{territory} (>0.85, asentado),
\texttt{moyo} (0.15--0.85, disputado), \texttt{neutral} (<0.15) --- pero
hasta ahora \texttt{filter\_moyos()} descartaba \texttt{territory} y
\texttt{neutral}. Se habilit\'o \texttt{territory} como segunda
categor\'ia (\texttt{filter\_territory()}, mismo mecanismo, sin volver
a dise\~nar nada) y se reconstruy\'o el cach\'e
(\texttt{cache\_full20\_cats}: 120 posiciones, 872 moyos, 312
territorios).

\subsection{Territorio asentado: la categor\'ia "f\'acil"}

\medskip
\begin{center}
\rowcolors{2}{white}{cTRalt}
\begin{tabular}{llrrr}
  \TH{Hamiltoniano} & \TH{Categor\'ia} & \TH{$R^2$ geom. sola} & \TH{$R^2$ completo} & \TH{$\Delta R^2$} \\[1pt]
  $H_{M1}$      & Moyo       & 0.312 & 0.589 & \good{0.277} \\
  $H_{M1}$      & Territorio & \good{0.670} & 0.789 & 0.119 \\
  $H_{OPT\_A}$  & Moyo       & 0.312 & 0.596 & \good{0.284} \\
  $H_{OPT\_A}$  & Territorio & \good{0.670} & 0.777 & 0.107 \\
\end{tabular}
\end{center}
\medskip

\begin{tcolorbox}[quotebox]
  Confirma la intuici\'on conceptual: en territorio asentado, la sola
  \textbf{geometr\'ia} (distancia a piedras/borde) ya explica el
  67\,\% de la varianza --- casi tautol\'ogico, un punto lejos de
  piedras enemigas y cerca de las propias casi siempre termina siendo
  territorio propio. El Hamiltoniano aporta bastante menos ah\'i
  ($\Delta R^2\approx0.11$--$0.12$) que en moyo ($\approx0.28$). El
  moyo sigue siendo, con evidencia, la categor\'ia \textbf{dif\'icil}
  donde el campo de $H$ realmente agrega algo que la geometr\'ia sola
  no puede --- territorio es m\'as un piso de control/sanidad que un
  reto real.
\end{tcolorbox}

\subsection{Fase de la partida (dentro del rango muestreado)}

El muestreo de posiciones ya restring\'ia deliberadamente a
jugadas entre el 15\,\% y el 90\,\% de cada partida
(\texttt{sample\_move\_numbers}, evita aperturas triviales y finales ya
resueltos) --- por lo que \textbf{no hay fuseki puro ni yose puro} en
los datos usados en esta secci\'on; solo un gradiente dentro de ese
rango medio. (La extensi\'on que s\'i muestrea fuera de ese rango ---
fuseki y joseki reales --- se presenta en la
Secci\'on~\ref{sec:joseki-fuseki}.) Segmentando el moyo en tercios de ese
rango:

\medskip
\begin{center}
\rowcolors{2}{white}{cTRalt}
\begin{tabular}{lrr}
  \TH{Fase (dentro de 15\%--90\%)} & \TH{$\Delta R^2$ $H_{M1}$} & \TH{$\Delta R^2$ $H_{OPT\_A}$} \\[1pt]
  Temprano (15--40\,\%) & \good{0.334} & \good{0.313} \\
  Medio (40--65\,\%)    & 0.152 & 0.154 \\
  Tard\'io (65--90\,\%) & \bad{0.102} & \bad{0.119} \\
\end{tabular}
\end{center}
\medskip

\begin{tcolorbox}[quotebox]
  El poder predictivo de \textbf{ambos} Hamiltonianos cae de forma
  mon\'otona conforme avanza la partida --- m\'as del triple entre
  temprano y tard\'io. Interpretaci\'on razonable: un moyo reci\'en
  formado depende m\'as de estructura simple (distancia, interacci\'on
  local), mientras que uno que sobrevive hasta tarde en la partida ya
  fue puesto a prueba t\'acticamente (invasiones, lecturas de vida-muerte)
  --- terreno donde un campo cl\'asico de 4 par\'ametros
  estructuralmente no puede competir.
\end{tcolorbox}

\subsection{Joseki y fuseki: jugadas realmente tempranas}
\label{sec:joseki-fuseki}

El muestreo de arriba exclu\'ia a prop\'osito el 3\,\%--15\,\%
inicial de cada partida. Para cubrir ese hueco se construy\'o
\texttt{early\_regions.py}: muestrea jugadas gen\'uinamente tempranas
(3\,\%--15\,\%) y detecta dos tipos de regi\'on --- \textbf{fuseki}
(mismo mecanismo de moyo/territorio de siempre, ownership-based, pero
en este momento temprano) y \textbf{joseki} (4 esquinas fijas del
tablero, geometr\'ia pura, independiente de \texttt{ownership}). Se
corri\'o sobre las 20 partidas (\texttt{cache\_early20}: 80 posiciones,
493 regiones fuseki, 320 regiones joseki) y se evalu\'o $\Delta R^2$
para los 47 Hamiltonianos del cat\'alogo (radio~9, mismo est\'andar que
el resto del informe):

\medskip
\begin{center}
\rowcolors{2}{white}{cTRalt}
\begin{tabular}{lrrr}
  \TH{Categor\'ia} & \TH{Momento} & \TH{$\Delta R^2$ $H_{M1}$} & \TH{$\Delta R^2$ $H_{OPT\_A}$} \\[1pt]
  Moyo (15--40\,\%, arriba)          & medio-temprano & 0.334 & 0.313 \\
  Fuseki (3--15\,\%, nuevo)          & muy temprano   & \good{0.339} & \good{0.327} \\
  Joseki (esquinas, 3--15\,\%, nuevo) & muy temprano  & 0.216 & \good{0.279} \\
\end{tabular}
\end{center}
\medskip

\begin{tcolorbox}[quotebox]
  \textbf{Confirma la tendencia de arriba, en vez de contradecirla}:
  fuseki (el momento a\'un m\'as temprano que ``temprano'') rinde igual
  o mejor que el tercio m\'as temprano ya medido, para ambos
  Hamiltonianos --- consistente con la ca\'ida mon\'otona de poder
  predictivo conforme avanza la partida. En joseki, $H_{OPT\_A}$ supera
  con m\'as margen a $H_{M1}$ (0.279 contra 0.216) que en cualquier
  otra categor\'ia del informe --- pero, siguiendo el est\'andar de la
  Secci\'on~\ref{sec:bootstrap}, esta diferencia \textbf{todav\'ia no
  tiene} un intervalo de confianza por partida; no se declara ganador
  hasta calcularlo.
\end{tcolorbox}

Como referencia de techo (no como candidato a adoptar): el mejor
Hamiltoniano aleatorio de todo el cat\'alogo, $H_{0115}$ --- ya
se\~nalado en la Secci\'on~3.3 como inestable entre radios --- alcanza
$\Delta R^2=0.449$ en fuseki y $0.471$ en joseki, los valores m\'as
altos medidos en todo el experimento en cualquier categor\'ia.

\subsection{Estado del pipeline: categor\'ias habilitadas hasta ahora}

\medskip
\begin{center}
\rowcolors{2}{white}{cTRalt}
\begin{tabular}{cll}
  \TH{\#} & \TH{Categor\'ia} & \TH{Estado} \\[1pt]
  1 & Moyo & \good{Habilitada} (Secciones~3--7) \\
  2 & Territorio asentado & \good{Habilitada} (Secci\'on~7.1) \\
  3 & Fase de partida (15\,\%--90\,\%) & \good{Habilitada} (Secci\'on~7.2) \\
  4 & Joseki (esquinas) & \good{Habilitada} (Secci\'on~\ref{sec:joseki-fuseki}) \\
  5 & Fuseki (apertura real, $<$15\,\%) & \good{Habilitada} (Secci\'on~\ref{sec:joseki-fuseki}) \\
  6 & Grupos de piedras / vida-muerte & pendiente (mecanismo nuevo) \\
  7 & Frontera activa / aji & pendiente (mecanismo nuevo) \\
  8 & Sente/gote (tempo) & \good{Habilitada} (Secci\'on~\ref{sec:sente-gote}, se\~nal d\'ebil) \\
\end{tabular}
\end{center}
\medskip

\subsection{Estado y siguientes pasos de la Parte~I}

\begin{enumerate}
  \item \good{Hecho}: Fase A termin\'o y encontr\'o $H_{OPT\_A}$, con
        mejor $\Delta R^2$ puntual que $H_{M1}$ en el cach\'e de 20
        partidas (Secci\'on~\ref{sec:optimizacion}).
  \item \bad{Corregido}: el bootstrap por partida
        (Secci\'on~\ref{sec:bootstrap}) muestra que esa ventaja
        \textbf{no es estad\'isticamente distinguible} del ruido de
        muestreo con solo 20 partidas --- los IC~95\,\% de $H_{M1}$ y
        $H_{OPT\_A}$ se traslapan casi por completo. $H_{OPT\_A}$ sigue
        siendo un candidato v\'alido y prometedor, no un reemplazo
        confirmado de $H_{M1}$.
  \item \good{Hecho}: se integr\'o \textbf{toda} la informaci\'on
        aprovechable de KataGo --- \texttt{ownershipStdev} (pondera la
        regresi\'on --- sube $\Delta R^2$ $\sim$30\,\% para ambos,
        \'unica se\~nal que aport\'o algo), \texttt{policy},
        \texttt{moveInfos}, \texttt{scoreStdev}, \texttt{winrate} y
        \texttt{scoreLead} (ninguna de estas \'ultimas cinco aport\'o
        se\~nal significativa) --- con \texttt{maxVisits} subido de
        250 a 600 (Secci\'on~\ref{sec:rich}).
  \item \good{Hecho}: se habilit\'o \texttt{territorio asentado} como
        segunda categor\'ia (Secci\'on~\ref{sec:categorias}) --- $H$
        aporta bastante menos ah\'i ($\Delta R^2\approx0.11$) que en
        moyo ($\approx0.28$), confirmando que moyo es la categor\'ia
        dif\'icil donde el modelo genuinamente compite. Tambi\'en se
        confirm\'o que el poder predictivo cae con la fase de la
        partida (0.33 temprano $\to$ 0.10 tard\'io, dentro del rango
        muestreado).
  \item \good{Hecho}: el \'optimo de Fase~A ($H_{OPT\_A}$) se integr\'o
        al visor interactivo existente como entrada del nuevo grupo
        \texttt{optimizado}, con su variedad 3D, puntos cr\'iticos
        marcados y ambos $\Delta R^2$ (r=1, r=9). Adem\'as, todos los
        47 Hamiltonianos del cat\'alogo se clasificaron por tipo de
        polinomio (tipo $H_{M1}$, espejo, nulo, cuadr\'atico puro,
        mixto) seg\'un el signo de $(\Sigma a,\Sigma b,b_{12},\Sigma c)$,
        con agrupaci\'on visual en el visor.
  \item Pendientes cruzados con el resto del pipeline (joseki, fuseki,
        sente/gote, m\'as partidas, Fase~B): ver el cierre unificado en
        la Secci\'on~\ref{sec:cierre}.
\end{enumerate}

% =================================================================================
\part{Por qu\'e el campo solo ve 4 de 7 par\'ametros: teor\'ia de invariantes}
\label{part:08}

% =================================================================================
\section{Motivaci\'on: de un hallazgo ad hoc a una herramienta sistem\'atica}

En la Parte~\ref{part:07} (Secci\'on~\ref{sec:degeneracion}) encontramos,
revisando el c\'odigo de \texttt{relaxation\_field}, que el campo que
predice moyos depende \'unicamente de 4 combinaciones de los 7
coeficientes de \texttt{cubic\_mixed}:
\[
\Sigma a = a_1+a_2, \qquad \Sigma b = b_{11}+b_{22}, \qquad b_{12},
\qquad \Sigma c = c_{112}+c_{122}.
\]
Ese hallazgo se hizo \textbf{a mano}: expandimos $H(s,q)+H(q,s)$
algebraicamente y vimos qu\'e sobreviv\'ia. Funciona, pero no explica
\emph{por qu\'e} son exactamente esas 4 combinaciones y no otras, ni
qu\'e pasar\'ia con una plantilla distinta. Esta parte deriva lo
mismo de manera \textbf{sistem\'atica y en un solo sentido}: primero
identificamos la \'unica simetr\'ia que hace falta para responder
"por qu\'e 4 y no 7" (posici\'on, $\sigma$), la resolvemos por
completo, y \textbf{despu\'es} a\~nadimos una segunda simetr\'ia
(color, $\tau$) para refinar la clasificaci\'on interna de esas 4 ---
sin volver a tocar el conteo principal.

% =================================================================================
\section{Los 7 polinomios de partida, uno por uno}
\label{sec:los7}

La plantilla \texttt{cubic\_mixed} es
\[
H(x,y) \;=\; a_1x + a_2y + b_{11}x^2 + b_{12}xy + b_{22}y^2
+ c_{112}x^2y + c_{122}xy^2,
\]
7 monomios, cada uno con su propio coeficiente libre:

\medskip
\begin{center}
\rowcolors{2}{white}{cTRalt}
\begin{tabular}{clc}
  \TH{\#} & \TH{Monomio} & \TH{Coeficiente} \\[1pt]
  1 & $x$     & $a_1$ \\
  2 & $y$     & $a_2$ \\
  3 & $x^2$   & $b_{11}$ \\
  4 & $xy$    & $b_{12}$ \\
  5 & $y^2$   & $b_{22}$ \\
  6 & $x^2y$  & $c_{112}$ \\
  7 & $xy^2$  & $c_{122}$ \\
\end{tabular}
\end{center}
\medskip

\'Estos son los 7 par\'ametros libres de cualquier Hamiltoniano de
esta familia --- el punto de partida completo. Las secciones que
siguen resuelven, en un solo camino hacia adelante, qu\'e les pasa a
estos 7 al pasar por \texttt{relaxation\_field}.

% =================================================================================
\section{¿Qu\'e es un grupo, aqu\'i?}

Un \textbf{grupo} es, para lo que necesitamos, una colecci\'on de
transformaciones que se pueden componer entre s\'i, donde:
\begin{itemize}
  \item hay una transformaci\'on "no hacer nada" (la \textbf{identidad} $e$);
  \item toda transformaci\'on tiene una \textbf{inversa} (deshacerla);
  \item componer transformaciones es \textbf{asociativo}.
\end{itemize}

Ejemplo m\'as simple: $\{+1,-1\}$ con la multiplicaci\'on usual es un
grupo de 2 elementos --- el m\'as peque\~no no trivial que existe. Es
exactamente el tama\~no de grupo que necesitamos primero.

% =================================================================================
\section{Paso 1: la simetr\'ia de posici\'on $\sigma$, y por qu\'e basta ella sola}
\label{sec:sigma}

\subsection{De d\'onde sale $\sigma$}

$H(x,y)$ punt\'ua la interacci\'on de un par de intersecciones
vecinas. Pero \textbf{no hay nada f\'isico que distinga cu\'al de las
dos es "la primera" ($i$) y cu\'al "la segunda" ($j$)} --- es una
elecci\'on de notaci\'on, no una propiedad del tablero. Por eso, en
\texttt{relaxation\_field} (\texttt{features.py:136}), cada
actualizaci\'on de un punto \textbf{siempre} usa la suma
$H(s,q)+H(q,s)$: ambos \'ordenes, sumados. Esa suma es, literalmente,
promediar sobre la transformaci\'on
\[
\sigma: (x,y) \;\longmapsto\; (y,x).
\]
Aplicarla dos veces regresa al original ($\sigma(\sigma(x,y))=(x,y)$)
--- $\sigma$ es una \textbf{involuci\'on}, y $\{e,\sigma\}$ es un
grupo de 2 elementos, el m\'as simple posible.

\subsection{Qu\'e le hace $\sigma$ a cada uno de los 7 monomios}

Un grupo de 2 elementos actuando sobre un espacio vectorial real
siempre lo separa en exactamente dos partes: la que \textbf{no cambia}
al aplicar la transformaci\'on (invariante) y la que \textbf{cambia de
signo} (anti-invariante) --- no hay una tercera opci\'on posible para
un grupo de orden 2. Aplicado monomio por monomio:

\medskip
\begin{center}
\rowcolors{2}{white}{cTRalt}
\begin{tabular}{lccc}
  \TH{Par de monomios} & \TH{Combinaci\'on $\sigma$-invariante} & \TH{Combinaci\'on $\sigma$-anti} & \TH{Total} \\[1pt]
  $\{x,y\}$       & $x{+}y$ (1)             & $x{-}y$ (1)     & 2 \\
  $\{x^2,y^2\}$   & $x^2{+}y^2$ (1)         & $x^2{-}y^2$ (1) & 2 \\
  $\{xy\}$        & $xy$ (1, ya sim\'etrico: $xy=yx$) & --- (0) & 1 \\
  $\{x^2y,xy^2\}$ & $x^2y{+}xy^2$ (1)       & $x^2y{-}xy^2$ (1) & 2 \\
  \hline
  \TH{Total} & \good{4} & \bad{3} & 7 \\
\end{tabular}
\end{center}
\medskip

El t\'ermino $xy$ es el \'unico sin pareja: como $xy=yx$ ya es
sim\'etrico por s\'i mismo, no genera una combinaci\'on
anti-invariante --- por eso el conteo final es 4 contra 3, no 3.5 y 3.5.

\subsection{Por qu\'e \texttt{relaxation\_field} solo ve la mitad invariante}

$H(s,q)+H(q,s)$ es, por construcci\'on, el doble de la proyecci\'on
sobre la parte $\sigma$-invariante: si $A(x,y)$ es cualquier
combinaci\'on \textbf{anti}-invariante (cambia de signo bajo $\sigma$),
entonces $A(s,q)+A(q,s) = A(s,q)-A(s,q) = 0$ exactamente --- se cancela
sin importar el valor de sus coeficientes. Verificado simb\'olicamente
con \texttt{sympy} sobre las tres combinaciones anti-invariantes de
\texttt{cubic\_mixed} ($x{-}y$, $x^2{-}y^2$, $x^2y{-}xy^2$, con
coeficientes libres $\Delta a=a_1{-}a_2$, $\Delta b=b_{11}{-}b_{22}$,
$\Delta c=c_{112}{-}c_{122}$): la suma da id\'enticamente cero.

\begin{tcolorbox}[quotebox]
  \textbf{Por qu\'e 4 y no 7, completo, con una sola simetr\'ia}: de
  los 7 monomios, 4 sobreviven al promedio $\sigma$ (combin\'andose en
  $\Sigma a=a_1{+}a_2$, $\Sigma b=b_{11}{+}b_{22}$, $b_{12}$ tal cual,
  y $\Sigma c=c_{112}{+}c_{122}$) y 3 se cancelan exactamente
  ($\Delta a$, $\Delta b$, $\Delta c$). Esta es la respuesta completa a
  "por qu\'e 4 y no 7" --- \textbf{no hace falta nada m\'as} que esta
  \'unica simetr\'ia de 2 elementos para derivarla.
\end{tcolorbox}

\subsection{Verificaci\'on emp\'irica: una predicci\'on falsable}

La teor\'ia predice algo m\'as fuerte que "solo importan 4
combinaciones": predice que un Hamiltoniano construido
\textbf{\'unicamente} con la parte anti-invariante ($\Delta a$,
$\Delta b$, o $\Delta c$ puros) deber\'ia dar $\Delta R^2=0$
\textbf{exacto}, sin importar la magnitud de sus coeficientes. Se
prob\'o sobre el cach\'e real de 20 partidas
(\texttt{cache\_full20\_cats}, radio~9):

\medskip
\begin{center}
\rowcolors{2}{white}{cTRalt}
\begin{tabular}{lr}
  \TH{Hamiltoniano} & \TH{$\Delta R^2$} \\[1pt]
  $\Delta a$ puro ($a_1{=}1,a_2{=}{-}1$, resto 0)         & \bad{0.000000} \\
  $\Delta b$ puro ($b_{11}{=}1,b_{22}{=}{-}1$, resto 0)   & \bad{0.000000} \\
  $\Delta c$ puro ($c_{112}{=}1,c_{122}{=}{-}1$, resto 0) & \bad{0.000000} \\
  Los tres combinados                                      & \bad{0.000000} \\
  \good{Control}: $\Sigma a$ puro ($a_1{=}1,a_2{=}1$)      & \good{0.000241} ($p=0.58$, distinto de 0) \\
\end{tabular}
\end{center}
\medskip

\begin{tcolorbox}[quotebox]
  Cero exacto, con precisi\'on de m\'aquina, para los cuatro
  Hamiltonianos puramente anti-invariantes --- no aproximadamente
  cero, \textbf{cero}. El control ($\Sigma a$ puro, que s\'i debe
  sobrevivir) da un valor peque\~no pero distinto de cero. La
  predicci\'on te\'orica se confirma sin ambig\"uedad sobre datos
  reales, y con esto qued\'o completamente resuelta la pregunta
  central de esta parte.
\end{tcolorbox}

% =================================================================================
\section{Paso 2: una segunda simetr\'ia, $\tau$, para refinar --- no para recontar}
\label{sec:tau}

Hasta aqu\'i, todo se deriv\'o con \textbf{una sola} simetr\'ia
($\sigma$). Esta secci\'on a\~nade una \textbf{segunda}, independiente,
para responder una pregunta distinta: de las 4 combinaciones que
sobrevivieron ($\Sigma a,\Sigma b,b_{12},\Sigma c$), ¿hay alguna
subestructura interna? ¿Por qu\'e $H_{M1}$ usa solo $\Sigma a,\Sigma c$
y Alvarado solo $b_{12}$?

\subsection{La simetr\'ia de color, $\tau$}

Tampoco hay nada f\'isico que distinga "negro" de "blanco" como
etiquetas --- son convenciones ($s=-1$ vs.\ $s=+1$). Invertir todos los
colores del tablero corresponde a
\[
\tau: (x,y) \;\longmapsto\; (-x,-y).
\]
Al igual que $\sigma$, es una involuci\'on ($\tau^2=e$).

\subsection{$\sigma$ y $\tau$ conmutan}

\[
\sigma(\tau(x,y)) = \sigma(-x,-y) = (-y,-x)
\qquad = \qquad
\tau(\sigma(x,y)) = \tau(y,x) = (-y,-x).
\]
Verificado tambi\'en sobre $H$ completo con \texttt{sympy}: aplicar
$\sigma$ despu\'es de $\tau$ da t\'ermino a t\'ermino lo mismo que
$\tau$ despu\'es de $\sigma$.

\subsection{Por qu\'e $\{e,\sigma,\tau,\sigma\tau\}$ es el grupo de Klein, no $\mathbb{Z}_4$}
\label{sec:klein}

Juntas, $\sigma$ y $\tau$ generan un grupo de 4 elementos,
$G=\{e,\sigma,\tau,\sigma\tau\}$. Hay \textbf{dos} grupos distintos de
4 elementos (salvo isomorfismo), y hay que distinguir cu\'al es \'este:

\begin{itemize}
  \item $\mathbb{Z}_4$ (c\'iclico): tiene un elemento de \textbf{orden 4}.
  \item $\mathbb{Z}_2\times\mathbb{Z}_2$ (\textbf{grupo de Klein}, $V_4$):
        \textbf{todos} sus elementos no triviales tienen orden 2.
\end{itemize}

Aqu\'i, $\sigma^2=e$, $\tau^2=e$, y $(\sigma\tau)^2=\sigma^2\tau^2=e$
(usando que conmutan) --- \textbf{los tres elementos no triviales
tienen orden 2}: es el grupo de Klein. Tabla de multiplicaci\'on
completa:

\medskip
\begin{center}
\rowcolors{2}{white}{cTRalt}
\begin{tabular}{c|cccc}
  $\cdot$ & $e$ & $\sigma$ & $\tau$ & $\sigma\tau$ \\
  \hline
  $e$          & $e$          & $\sigma$      & $\tau$        & $\sigma\tau$ \\
  $\sigma$     & $\sigma$     & $e$           & $\sigma\tau$  & $\tau$ \\
  $\tau$       & $\tau$       & $\sigma\tau$  & $e$           & $\sigma$ \\
  $\sigma\tau$ & $\sigma\tau$ & $\tau$        & $\sigma$      & $e$ \\
\end{tabular}
\end{center}
\medskip

Cada elemento aparece exactamente una vez por fila/columna (grupo
v\'alido) y la diagonal es puro $e$ --- la firma exacta del grupo de
Klein.

\begin{tcolorbox}[quotebox]
  \textbf{Por qu\'e importa que sea Klein y no c\'iclico}: en
  $\mathbb{Z}_4$ las representaciones irreducibles usan n\'umeros
  complejos (ra\'ices de la unidad) porque hay un elemento de orden 4.
  En Klein, como todo tiene orden $\leq2$, todas las representaciones
  son \textbf{reales}, con valores $\pm1$ \'unicamente --- justo lo que
  necesitamos, porque nuestros coeficientes son n\'umeros reales.
\end{tcolorbox}

\subsection{Las 4 piezas de Klein, y d\'onde queda cada mon\'omio}

Como $G$ es abeliano de orden 4, tiene exactamente 4 representaciones
irreducibles, todas de dimensi\'on 1 --- cada una es un \textbf{car\'acter}
que asigna $\pm1$ a $\sigma$ y a $\tau$ por separado:

\medskip
\begin{center}
\rowcolors{2}{white}{cTRalt}
\begin{tabular}{lccl}
  \TH{Car\'acter} & \TH{$\sigma\mapsto$} & \TH{$\tau\mapsto$} & \TH{Significado} \\[1pt]
  $\chi_{++}$ & $+1$ & $+1$ & sim\'etrico en posici\'on y en color \\
  $\chi_{+-}$ & $+1$ & $-1$ & sim\'etrico en posici\'on, impar en color \\
  $\chi_{-+}$ & $-1$ & $+1$ & antisim\'etrico en posici\'on, par en color \\
  $\chi_{--}$ & $-1$ & $-1$ & antisim\'etrico en ambas \\
\end{tabular}
\end{center}
\medskip

El operador de Reynolds $P_\chi(H)=\frac{1}{4}\sum_{g\in G}\chi(g)\,(g{\cdot}H)$
proyecta sobre cada pieza. Aplic\'andolo a \texttt{cubic\_mixed}
(verificado simb\'olicamente, suma exacta $=H$):

\begin{tcolorbox}[codebox]
$P_{++}(H) = \dfrac{\Sigma b}{2}(x^2{+}y^2) + b_{12}\,xy$
\hfill \textit{(2 dim.)} \\[4pt]
$P_{+-}(H) = \dfrac{\Sigma a}{2}(x{+}y) + \dfrac{\Sigma c}{2}(x^2y{+}xy^2)$
\hfill \textit{(2 dim.)} \\[4pt]
$P_{-+}(H) = \dfrac{\Delta b}{2}(x^2{-}y^2)$
\hfill \textit{(1 dim.)} \\[4pt]
$P_{--}(H) = \dfrac{\Delta a}{2}(x{-}y) + \dfrac{\Delta c}{2}(x^2y{-}xy^2)$
\hfill \textit{(2 dim.)}
\end{tcolorbox}

N\'otese: $P_{++}\oplus P_{+-}$ (dimensi\'on $2+2=4$) es exactamente la
parte $\sigma$-invariante de la Secci\'on~\ref{sec:sigma} --- $\tau$ no
cambi\'o cu\'antas dimensiones sobreviven, solo las \textbf{parti\'o en
dos} seg\'un el color. Y $P_{-+}\oplus P_{--}$ ($1+2=3$) es la parte
$\sigma$-anti que ya sab\'iamos que muere, ahora tambi\'en partida en
dos por color (irrelevante para \texttt{relaxation\_field}, pero
consistente).

% =================================================================================
\section{Tabla final: los 7 monomios, de principio a fin}

Uniendo ambos pasos, as\'i queda cada uno de los 7 monomios
originales:

\medskip
\begin{center}
\rowcolors{2}{white}{cTRalt}
\begin{tabular}{clcccl}
  \TH{\#} & \TH{Monomio} & \TH{Coef.} & \TH{Combinaci\'on que sobrevive} & \TH{Pieza de Klein} & \TH{¿Sobrevive?} \\[1pt]
  1 & $x$    & $a_1$     & $\Sigma a = a_1{+}a_2$         & $P_{+-}$ & \good{S\'i} \\
  2 & $y$    & $a_2$     & $\Sigma a = a_1{+}a_2$         & $P_{+-}$ & \good{S\'i} \\
  3 & $x^2$  & $b_{11}$  & $\Sigma b = b_{11}{+}b_{22}$   & $P_{++}$ & \good{S\'i} \\
  4 & $xy$   & $b_{12}$  & $b_{12}$ (sin pareja, intacto) & $P_{++}$ & \good{S\'i} \\
  5 & $y^2$  & $b_{22}$  & $\Sigma b = b_{11}{+}b_{22}$   & $P_{++}$ & \good{S\'i} \\
  6 & $x^2y$ & $c_{112}$ & $\Sigma c = c_{112}{+}c_{122}$ & $P_{+-}$ & \good{S\'i} \\
  7 & $xy^2$ & $c_{122}$ & $\Sigma c = c_{112}{+}c_{122}$ & $P_{+-}$ & \good{S\'i} \\
\end{tabular}
\end{center}
\medskip

\begin{tcolorbox}[defbox]
\textbf{Las 4 combinaciones que sobreviven, tal cual:}
\begin{center}
\rowcolors{2}{white}{cTRalt}
\begin{tabular}{lll}
  \TH{Combinaci\'on} & \TH{F\'ormula} & \TH{Viene de} \\[1pt]
  $\Sigma a$ & $a_1+a_2$       & coeficientes de $x$ y $y$ \\
  $\Sigma b$ & $b_{11}+b_{22}$ & coeficientes de $x^2$ y $y^2$ \\
  $b_{12}$   & $b_{12}$        & coeficiente de $xy$ \\
  $\Sigma c$ & $c_{112}+c_{122}$ & coeficientes de $x^2y$ y $xy^2$ \\
\end{tabular}
\end{center}
\end{tcolorbox}

\begin{center}
\rowcolors{2}{white}{cTRalt}
\begin{tabular}{lll}
  \TH{Las 3 que mueren} & \TH{F\'ormula} & \TH{Viene de} \\[1pt]
  $\Delta a$ & $a_1-a_2$       & coeficientes de $x$ y $y$ \\
  $\Delta b$ & $b_{11}-b_{22}$ & coeficientes de $x^2$ y $y^2$ \\
  $\Delta c$ & $c_{112}-c_{122}$ & coeficientes de $x^2y$ y $xy^2$ \\
\end{tabular}
\end{center}

% =================================================================================
\section{Lo que $\tau$ explica que $\sigma$ sola no explicaba}
\label{sec:paraque2}

\subsection{Correcci\'on importante: $H_{M1}$ usa \textbf{dos} particiones, no una}

Una primera lectura sugiere que, como \texttt{sparse\_cubic}/$H_{M1}$
solo usa $a_1,a_2,c_{112},c_{122}$ (nunca $b_{11},b_{12},b_{22}$), su
plantilla vivir\'ia \'integramente en $P_{+-}$. \textbf{Eso es
incorrecto} si se toma $H_{M1}$ con sus coeficientes reales
--- hay que revisar las \textbf{dos} particiones $\tau$-impares
($P_{+-}$ \emph{y} $P_{--}$), no solo una. Descomponiendo
$H_{M1}=x+2y-x^2y-xy^2$ (verificado con \texttt{sympy}):

\begin{tcolorbox}[codebox]
$P_{+-}(H_{M1}) = \dfrac{3}{2}(x{+}y) \,-\, 1\cdot(x^2y{+}xy^2)$
\hfill ($\Sigma a{=}3,\ \Sigma c{=}{-}2$) \\[4pt]
$P_{--}(H_{M1}) = -\dfrac{1}{2}(x{-}y) \,+\, 0\cdot(x^2y{-}xy^2)$
\hfill ($\Delta a{=}{-}1,\ \Delta c{=}0$) \\[4pt]
$P_{++}(H_{M1}) = 0, \qquad P_{-+}(H_{M1}) = 0$
\end{tcolorbox}

$H_{M1} = P_{+-}(H_{M1}) + P_{--}(H_{M1})$, y \textbf{ninguna de las
dos partes es cero}: $a_1=1\neq a_2=2$, as\'i que $\Delta a=a_1-a_2=-1$
no se anula. $H_{M1}$ s\'i es completamente $\tau$-\textbf{impar}
(vive en la uni\'on $P_{+-}\oplus P_{--}$, nunca toca $P_{++}$ ni
$P_{-+}$ --- eso s\'i es correcto, y por eso $H(-x,-y)=-H(x,y)$ se
cumple exactamente, verificado). Pero \textbf{dentro} de ese lado
impar-en-color, $H_{M1}$ se reparte entre las dos particiones de
$\sigma$: la mayor parte en $P_{+-}$ (lo que \texttt{relaxation\_field}
s\'i ve), y una parte menor pero no nula en $P_{--}$ (lo que
\texttt{relaxation\_field} \textbf{nunca} ve --- se cancela en
$H(s,q)+H(q,s)$ sin importar que $H_{M1}$, como polinomio completo,
s\'i la tenga).

\begin{tcolorbox}[quotebox]
  \textbf{Esto no es un caso aislado}: se verific\'o que
  $H_{OPT\_A}$ (el \'optimo de Fase~A) \textbf{tambi\'en} tiene
  componente $P_{--}$ no nula ($\Delta a\approx-0.340$,
  $\Delta c\approx1.105$). Ning\'un Hamiltoniano probado hasta ahora en
  este experimento vive puramente en $P_{+-}$ --- la clasificaci\'on
  por "tipo" (Secci\'on siguiente) describe correctamente el signo de
  $\Sigma a,\Sigma c$ (lo que importa para $\Delta R^2$), pero
  \textbf{no} implica que el resto de $H$ sea cero; solo que es
  invisible para \texttt{relaxation\_field}.
\end{tcolorbox}

\subsection{¿Cu\'ando es puro y cu\'ando mixto? La regla general}
\label{sec:regla-mixto}

Lo anterior no es una peculiaridad de $H_{M1}$: es la regla para
\textbf{casi cualquier} Hamiltoniano. Cada par de coeficientes
originales ($a_1,a_2$), ($b_{11},b_{22}$), ($c_{112},c_{122}$) se
reparte en su suma ($\Sigma$, sobrevive) y su resta ($\Delta$, muere)
--- un Hamiltoniano es \textbf{puro} en una sola pieza de Klein
\'unicamente cuando todas sus $\Delta$ correspondientes son cero
exactamente:

\medskip
\begin{center}
\rowcolors{2}{white}{cTRalt}
\begin{tabular}{lll}
  \TH{Caso} & \TH{¿Por qu\'e es puro?} & \TH{Pieza} \\[1pt]
  Alvarado ($b_{12}=1$, resto 0) & $a_1{=}a_2{=}0$, $b_{11}{=}b_{22}{=}0$: & $P_{++}$ puro \\
   & todos los $\Delta$ son 0 trivialmente & \\
  Clase unitaria $(+1,+1,-1,-1)$ & $a_1{=}a_2{=}1\Rightarrow\Delta a{=}0$; & $P_{+-}$ puro \\
  (signos de $H_{M1}$, del barrido) & $c_{112}{=}c_{122}{=}{-}1\Rightarrow\Delta c{=}0$ & \\
  $H_{M1}$ real ($a_1{=}1,a_2{=}2$) & $a_1\neq a_2\Rightarrow\Delta a{=}{-}1\neq0$ & $P_{+-}\oplus P_{--}$, mixto \\
  $H_{OPT\_A}$ (optimizado) & $a_1\neq a_2$ y $c_{112}\neq c_{122}$ & $P_{+-}\oplus P_{--}$, mixto \\
  Cualquier candidato aleatorio & coeficientes continuos: la & casi seguro mixto \\
  (los 44 + 25 de este experimento) & probabilidad de $a_1{=}a_2$ exacto es 0 & \\
\end{tabular}
\end{center}
\medskip

\begin{tcolorbox}[quotebox]
  \textbf{Conclusi\'on pr\'actica}: la pureza (vivir en una sola pieza
  de Klein) es la \textbf{excepci\'on}, no la norma --- ocurre solo
  cuando los coeficientes se eligieron a prop\'osito con esa simetr\'ia
  extra (como Alvarado, o las clases unitarias del barrido con
  coeficientes iguales por pareja). $H_{M1}$, $H_{OPT\_A}$, y
  virtualmente todos los candidatos aleatorios de \texttt{cubic\_mixed}
  y \texttt{sparse\_cubic} son mixtos: tienen una parte que
  \texttt{relaxation\_field} ve (la proyecci\'on $\sigma$-invariante) y
  una parte que nunca ve, pero que sigue siendo parte real del
  polinomio.
\end{tcolorbox}

\subsubsection{Conteo exacto sobre los 47 Hamiltonianos catalogados}

Se calcul\'o $\Delta a=a_1{-}a_2$, $\Delta b=b_{11}{-}b_{22}$,
$\Delta c=c_{112}{-}c_{122}$ para \textbf{cada uno} de los 47
Hamiltonianos del cat\'alogo del visor (los 44 \texttt{cubic\_mixed} +
$H_{M1}$ + Alvarado + $H_{OPT\_A}$), no solo para los 3 ejemplos
citados arriba:

\medskip
\begin{center}
\rowcolors{2}{white}{cTRalt}
\begin{tabular}{lr}
  \TH{Categor\'ia} & \TH{Cantidad (de 47)} \\[1pt]
  \good{Puros} (los 3 $\Delta$ son exactamente 0) & \good{1} \\
  \bad{Mixtos} (al menos un $\Delta\neq0$) & \bad{46} \\
\end{tabular}
\end{center}
\medskip

\begin{tcolorbox}[quotebox]
  El \'unico Hamiltoniano puro del cat\'alogo completo es
  \textbf{Alvarado} ($H_{0001}$) --- y solo porque su plantilla
  \texttt{quadratic} fija $a_1=a_2=0$ y $b_{11}=b_{22}=0$
  trivialmente, no por dise\~no expl\'icito. Los otros 46 --- incluido
  $H_{M1}$ --- son mixtos: tienen una parte $\sigma$-invariante (la que
  $\Delta R^2$ mide) y una parte $\sigma$-anti que \texttt{relaxation\_field}
  nunca ve, pero que existe como parte real del polinomio. Ser
  "mixto" no es la excepci\'on entre los candidatos aleatorios --- es
  virtualmente universal.
\end{tcolorbox}

\textbf{Consecuencia para la clasificaci\'on por "tipo" del visor}: las
etiquetas ("tipo $H_{M1}$", "tipo espejo", \ldots) se calculan
\'unicamente a partir del signo de $\Sigma a,\Sigma b,b_{12},\Sigma c$
--- nunca miran $\Delta a,\Delta b,\Delta c$. Dos Hamiltonianos con el
mismo signo de $\Sigma a,\Sigma c$ pero distinto $\Delta a,\Delta c$
reciben la \textbf{misma} etiqueta de tipo, aunque sean polinomios
distintos como objetos matem\'aticos completos. La etiqueta describe
correctamente la parte relevante para $\Delta R^2$; no pretende
describir el Hamiltoniano completo.

\subsection{Qu\'e s\'i explica $\tau$, entonces}

$\tau$ divide el espacio en dos mitades por color (par/impar), y
$\sigma$ en dos mitades por posici\'on (sim\'etrico/anti) ---
\textbf{independientes entre s\'i}. $P_{+-}$ y $P_{--}$ juntas
($\Sigma a,\Sigma c,\Delta a,\Delta c$) son exactamente la mitad
$\tau$-impar --- la \'unica que puede usar \texttt{sparse\_cubic}, ya
que su plantilla no incluye t\'erminos $b$. Que $H_{M1}$ cargue
tambi\'en una parte $\Delta a$ no nula no rompe esa pertenencia a la
mitad impar; solo confirma que, \textbf{dentro} de esa mitad, $H_{M1}$
no est\'a alineado exactamente con el eje $\sigma$-invariante. $P_{++}$
(par en color) contiene $\Sigma b$ y $b_{12}$ --- Alvarado ($H=xy$)
vive aqu\'i, y en ese caso s\'i por completo, porque $b_{11}=b_{22}=0$
hace que $\Delta b=0$ trivialmente.

La clasificaci\'on por "tipo de polinomio" ya construida en el visor
(tipo $H_{M1}$, tipo espejo, cuadr\'atico puro, nulo) describe
correctamente el signo de $(\Sigma a,\Sigma b,b_{12},\Sigma c)$ ---
la parte que importa para $\Delta R^2$ --- pero, a la luz de esta
secci\'on, debe leerse como "la proyecci\'on $\sigma$-invariante tiene
este signo", no como "el Hamiltoniano completo es esto".

% =================================================================================
\section{Implicaci\'on pr\'actica y alcance}

\textbf{Para Fase~B}: buscar directamente en las 4 dimensiones
efectivas ($\Sigma a,\Sigma b,b_{12},\Sigma c$) en vez de las 7 crudas
es estrictamente m\'as eficiente, con garant\'ia te\'orica (no solo
emp\'irica) de que no se pierde nada.

\textbf{Generalidad de la herramienta}: el mismo procedimiento ---
identificar las simetr\'ias del problema, reconocer el grupo que
generan, descomponer en representaciones irreducibles, proyectar ---
aplica a cualquier plantilla de Hamiltoniano futura, no solo
\texttt{cubic\_mixed}.

\begin{tcolorbox}[quotebox]
  \textbf{Advertencia de alcance (igual que en la Parte~\ref{part:07})}:
  esta reducci\'on aplica \'unicamente al pipeline de predicci\'on de
  moyo (\texttt{relaxation\_field}). El an\'alisis topol\'ogico de
  $\Gamma(H)$ (nodos $A_1$, puntos cr\'iticos) punt\'ua $H(x,y)$
  directamente, sin promediar sobre $\sigma$ --- ah\'i s\'i importan
  los 7 coeficientes por separado.
\end{tcolorbox}

% =================================================================================
\section{Cierre general y pr\'oximos pasos}
\label{sec:cierre}

\subsection{Sente/gote (tempo): resultados}
\label{sec:sente-gote}

A diferencia de moyo/territorio/fase/joseki/fuseki --- todas
categor\'ias \textbf{espaciales} (una regi\'on del tablero) --- sente/gote
se defini\'o como una se\~nal a \textbf{nivel de posici\'on completa}.
Se confirm\'o emp\'iricamente que \texttt{rootInfo.scoreLead} de KataGo
ya viene en perspectiva fija por partida (no relativa al jugador en
turno: se verific\'o que no cambia de signo cuando \texttt{currentPlayer}
alterna entre B y W en jugadas consecutivas), lo que permite calcular
con seguridad el cambio de marcador entre momentos muestreados
consecutivos --- combinando \texttt{cache\_early20} y
\texttt{cache\_full20\_cats} para cubrir toda la partida (140 intervalos
en 20 partidas) --- como una tasa de \textbf{volatilidad de tempo}
($|\Delta\texttt{scoreLead}|/\Delta\texttt{jugadas}$). Se prob\'o si la
varianza del campo de cada Hamiltoniano sobre \textbf{todo el tablero}
(no una regi\'on) predice esa volatilidad, con el mismo bootstrap por
partida de la Secci\'on~\ref{sec:bootstrap}:

\medskip
\begin{center}
\rowcolors{2}{white}{cTRalt}
\begin{tabular}{lrr}
  \TH{Hamiltoniano} & \TH{$\Delta R^2$} & \TH{IC 95\,\% (bootstrap por partida)} \\[1pt]
  $H_{M1}$                        & 0.042 & $[0.004,\ 0.111]$ \\
  $H_{OPT\_A}$                    & 0.035 & $[0.000,\ 0.168]$ \\
  Alvarado                        & \bad{0.002} & $[0.000,\ 0.016]$ \\
  Mejor de 47 ($H_{0140}$)        & 0.062 & $[0.007,\ 0.144]$ \\
\end{tabular}
\end{center}
\medskip

\begin{tcolorbox}[quotebox]
  \textbf{Se\~nal real pero muy d\'ebil} --- un orden de magnitud menor
  que en cualquier otra categor\'ia del informe (moyo/fuseki/joseki
  rondan $0.2$--$0.4$; aqu\'i el mejor de los 47 Hamiltonianos apenas
  llega a $0.06$). El IC~95\,\% de $H_{M1}$ no toca cero
  ($[0.004,0.111]$), as\'i que hay ind\'icio de se\~nal real, pero
  $H_{OPT\_A}$ s\'i lo toca casi ($[0.000,0.168]$) y Alvarado es
  esencialmente nulo. Interpretaci\'on honesta: la varianza del campo
  sobre todo el tablero contiene \textbf{algo} de informaci\'on sobre
  cu\'anto va a moverse el marcador pronto, pero mucho menos de lo que
  la misma familia de Hamiltonianos captura sobre d\'onde va a quedar
  el territorio. De las 8 categor\'ias de la taxonom\'ia de Go
  contempladas en este proyecto, sente/gote es, hasta ahora, la m\'as
  d\'ebil.
\end{tcolorbox}

\subsection{Fase~B: optimizando directamente en las 4 dimensiones efectivas}
\label{sec:faseb}

Aprovechando la garant\'ia te\'orica de la Parte~\ref{part:08} --- que
$\Delta R^2$ depende \'unicamente de $(\Sigma a,\Sigma b,b_{12},\Sigma c)$,
nunca de los 7 coeficientes crudos por separado --- se lanz\'o una
segunda ronda de optimizaci\'on directa, reparametrizada en esas 4
dimensiones, sobre la familia \texttt{cubic\_mixed} \textbf{completa}
(a diferencia de Fase~A, que solo explor\'o \texttt{sparse\_cubic}, sin
los t\'erminos $b_{11},b_{12},b_{22}$). Antes de lanzar la b\'usqueda se
verific\'o la reparametrizaci\'on: evaluar $H_{M1}$ mapeado a
$(\Sigma a,\Sigma b,b_{12},\Sigma c)=(3,0,0,-2)$ reproduce
$\Delta R^2=0.3152$, id\'entico al valor medido con sus coeficientes
crudos originales. Mismo protocolo que Fase~A: b\'usqueda sembrada con
$H_{M1}$, $H_{OPT\_A}$ y Alvarado sobre el cach\'e chico de 6 partidas,
validaci\'on autom\'atica sobre las 20 partidas completas al terminar.

\textbf{Resultado} (692 evaluaciones, $\sim$102 min, sobre el cach\'e de
6 partidas):

\medskip
\begin{center}
\rowcolors{2}{white}{cTRalt}
\begin{tabular}{ll}
  \TH{Combinaci\'on efectiva} & \TH{Valor \'optimo} \\[1pt]
  $\Sigma a$ & $-0.2483$ \\
  $\Sigma b$ & $\phantom{-}3.7286$ \\
  $b_{12}$   & $-3.0000$ \small (\bad{en el borde de su rango} $[-3,3]$) \\
  $\Sigma c$ & $-0.9660$ \\
  $\Delta R^2$ (cach\'e de 6 partidas) & \good{0.5477} \\
\end{tabular}
\end{center}
\medskip

Muy superior al $0.3523$ que hab\'ia encontrado Fase~A sobre el mismo
cach\'e --- la primera confirmaci\'on de que explorar el eje
$\Sigma b, b_{12}$ (los t\'erminos cuadr\'aticos, invisibles para
\texttt{sparse\_cubic}) s\'i aporta algo que $H_{M1}$/$H_{OPT\_A}$ no
ten\'ian.

\subsubsection{Validaci\'on sobre las 20 partidas completas, con bootstrap por partida}

Llamemos $H_{OPT\_B}$ al Hamiltoniano resultante (coeficientes crudos:
$a_1=a_2=-0.1241$, $b_{11}=b_{22}=1.8643$, $b_{12}=-3.0$,
$c_{112}=c_{122}=-0.4830$). Aplicando el mismo est\'andar de rigor de
la Secci\'on~\ref{sec:barrido9} (validar siempre fuera del subset donde
se encontr\'o) y de la Secci\'on~\ref{sec:bootstrap} (bootstrap por
partida, no por fila):

\medskip
\begin{center}
\rowcolors{2}{white}{cTRalt}
\begin{tabular}{lrr}
  \TH{Hamiltoniano} & \TH{$\Delta R^2$ (20 partidas, puntual)} & \TH{IC 95\,\% (bootstrap por partida)} \\[1pt]
  $H_{M1}$      & 0.298 & $[0.223,\ 0.358]$ \\
  $H_{OPT\_A}$  & 0.317 & $[0.227,\ 0.380]$ \\
  $H_{OPT\_B}$  & \good{0.421} & \good{$[0.339,\ 0.486]$} \\
\end{tabular}
\end{center}
\medskip

\begin{tcolorbox}[quotebox]
  \textbf{A diferencia de $H_{OPT\_A}$ contra $H_{M1}$, aqu\'i el
  traslape es peque\~no, no casi total}: el l\'imite inferior del
  IC~95\,\% de $H_{OPT\_B}$ ($0.339$) queda por encima del punto
  estimado de $H_{M1}$ ($0.298$) y muy cerca del l\'imite superior de
  su intervalo ($0.358$) --- la zona de traslape real con $H_{M1}$ es
  angosta ($[0.339,0.358]$), y con $H_{OPT\_A}$ tambi\'en
  ($[0.339,0.380]$). Es la evidencia \textbf{m\'as fuerte de todo el
  experimento} de que salir del subespacio \texttt{sparse\_cubic} (sin
  t\'erminos $b$) hacia la familia completa \texttt{cubic\_mixed}
  produce una mejora real, no solo un punto estimado favorable por
  ruido de muestreo.
\end{tcolorbox}

\subsubsection{Aclaraci\'on: $\Delta R^2$ no es el $R^2$ total}
\label{sec:aclaracion-r2}

Todo $\Delta R^2$ de este informe mide \'unicamente lo que el
Hamiltoniano \textbf{agrega} sobre el modelo base de geometr\'ia, no el
poder predictivo total del modelo completo. Para $H_{OPT\_B}$ en moyo,
$R^2_{\text{geometr\'ia sola}}=0.312$; sumando el Hamiltoniano,

\[
R^2_{\text{completo}} = R^2_{\text{geometr\'ia}} + \Delta R^2
= 0.312 + 0.421 = \good{0.733},
\qquad
r=\sqrt{R^2_{\text{completo}}} \approx \good{0.86}.
\]

\begin{tcolorbox}[quotebox]
  Le\'ido de forma aislada, $\Delta R^2=0.42$ puede parecer un n\'umero
  bajo en una escala de 0 a 1 --- pero no es la m\'etrica de "qu\'e tan
  bien predice el modelo completo", sino de "cu\'anto aporta el
  Hamiltoniano por encima de la geometr\'ia". El $R^2$ del modelo
  completo (geometr\'ia + $H_{OPT\_B}$) es $0.733$, equivalente a una
  correlaci\'on $r\approx0.86$ entre la predicci\'on y el territorio
  real medido por KataGo --- y el incremento en s\'i tiene
  $p=1.1\times10^{-16}$ (prueba $F$) con IC~95\,\% $[0.339,0.486]$
  (bootstrap por partida) que no se acerca a cero en ning\'un extremo:
  evidencia estad\'istica fuerte de correlaci\'on real, no ausencia de
  ella. El objetivo del experimento nunca fue igualar a KataGo (que
  hace lectura t\'actica completa), sino verificar si un Hamiltoniano
  cl\'asico de Ising captura algo real del territorio m\'as all\'a de
  la geometr\'ia --- y lo hace, de forma robusta.
\end{tcolorbox}

\begin{tcolorbox}[pendingbox]
  \textbf{Advertencia metodol\'ogica}: $b_{12}=-3.0$ cay\'o
  \textbf{exactamente} en el borde de su rango permitido ($[-3,3]$,
  el mismo rango usado para generar los 44 candidatos
  \texttt{cubic\_mixed} del cat\'alogo, por consistencia). Esto es
  ind\'icio de que el \'optimo verdadero (sin esa cota) podr\'ia ser
  a\'un mejor --- pero ampliar el rango de $b_{12}$ m\'as all\'a de lo
  que se usa para el resto del cat\'alogo rompe la comparabilidad
  directa entre $H_{OPT\_B}$ y los otros 47 Hamiltonianos. Queda como
  decisi\'on expl\'icita de dise\~no, no como b\'usqueda adicional
  autom\'atica: ¿ampliar el rango y aceptar que $H_{OPT\_B}$ ya no es
  directamente comparable al cat\'alogo, o mantener el rango y aceptar
  que este resultado es el mejor \textbf{dentro} de la familia
  estudiada, no un \'optimo global sin restricciones?
\end{tcolorbox}

\subsection{Estado actual, en una frase}

\good{Hecho}: Fase~B (Secci\'on~\ref{sec:faseb}) encontr\'o $H_{OPT\_B}$,
con $\Delta R^2=0.421$ sobre las 20 partidas e IC~95\,\%
$[0.339,0.486]$ --- por encima, con traslape angosto, de $H_{M1}$ y
$H_{OPT\_A}$. Es la evidencia m\'as fuerte del experimento de que
ampliar a la familia completa \texttt{cubic\_mixed} (con t\'erminos
$b$) ayuda. Todo lo dem\'as de esta secci\'on son \textbf{decisiones
pendientes}, no tareas ya definidas --- se presentan como caminos
alternativos, no como una lista \'unica ordenada, porque elegir entre
ellos depende de qu\'e pregunta interesa m\'as responder despu\'es.

% =================================================================================
\section{Caminos posibles a seguir}
\label{sec:caminos}

\subsection{Camino 1 --- Resolver el borde de $b_{12}$ en Fase B}

$b_{12}=-3.0$ cay\'o exactamente en el l\'imite de su rango permitido
(Secci\'on~\ref{sec:faseb}). Dos opciones, mutuamente excluyentes:

\begin{itemize}
  \item \textbf{Ampliar el rango} (p.ej. a $[-6,6]$, igual que
        $\Sigma b$) y repetir Fase~B --- costo: otra corrida de
        $\sim$100 min, y $H_{OPT\_B}$ dejar\'ia de ser directamente
        comparable en $b_{12}$ con los otros 47 Hamiltonianos del
        cat\'alogo (que s\'i respetan $[-3,3]$).
  \item \textbf{Mantener el rango actual} y aceptar $H_{OPT\_B}$ como
        el mejor candidato \emph{dentro} de la familia estudiada, con
        comparabilidad limpia frente al resto del cat\'alogo, sin saber
        si existe un \'optimo a\'un mejor m\'as all\'a del l\'imite.
\end{itemize}

\subsection{Camino 2 --- Rehacer el an\'alisis de variedades $\Gamma(H)$, pero sobre el polinomio reducido}
\label{sec:camino-variedades}

El an\'alisis topol\'ogico del experimento~06 (nodos $A_1$, puntos
cr\'iticos, fibra de Milnor de $\Gamma(H)$) se calcul\'o siempre sobre
el Hamiltoniano \textbf{crudo de 7 par\'ametros}, y no correlacion\'o
con $\Delta R^2$ real --- el grupo "Frente~1", elegido justo por ese
criterio, tuvo el \textbf{peor} desempe\~no emp\'irico
(Secci\'on~3.1). La Parte~\ref{part:08} explica \textbf{por qu\'e} ten\'ia
que pasar eso: \texttt{relaxation\_field} solo ve la proyecci\'on
$\sigma$-invariante de 4 dimensiones ($\Sigma a,\Sigma b,b_{12},\Sigma c$);
el an\'alisis de $\Gamma(H)$, al no promediar sobre $\sigma$, mezcla esas
4 dimensiones relevantes con las 3 ($\Delta a,\Delta b,\Delta c$) que son
\textbf{ruido puro} para la predicci\'on. Cualquier clasificaci\'on
topol\'ogica del Hamiltoniano completo est\'a, por construcci\'on,
contaminada por dimensiones invisibles al predictor real.

\begin{tcolorbox}[quotebox]
  \textbf{Camino nuevo, no probado a\'un}: repetir el mismo an\'alisis
  de variedades (puntos cr\'iticos, tipo de nodo, fibra de Milnor), no
  sobre $H(x,y)$ crudo, sino sobre el \textbf{polinomio reducido}
  $\Sigma a(x{+}y)+\Sigma b(x^2{+}y^2)+b_{12}xy+\Sigma c(x^2y{+}xy^2)$
  --- exactamente el objeto que \texttt{relaxation\_field} eval\'ua, sin
  las 3 dimensiones invisibles. Ser\'ia la primera vez que se prueba si
  la topolog\'ia correlaciona con $\Delta R^2$ usando el objeto correcto
  (antes ni siquiera se sab\'ia qu\'e dimensiones filtrar). Podr\'ia
  explicar, por ejemplo, por qu\'e $H_{OPT\_B}$ (con $b_{12}$ en el
  extremo de su rango) rinde tanto mejor --- quiz\'a su variedad
  reducida tiene una estructura de puntos cr\'iticos cualitativamente
  distinta a la de $H_{M1}$/$H_{OPT\_A}$.
\end{tcolorbox}

\textbf{Costo}: moderado --- reutiliza el c\'odigo de puntos
cr\'iticos/Milnor del experimento~06, pero hay que rederivar $\Gamma$
sobre el polinomio reducido de 4 par\'ametros en vez del de 7, y
recalcularlo para el cat\'alogo completo y para el espacio que explora
Fase~B. Puede volver a salir nulo --- pero ser\'ia la primera prueba
limpia de esta hip\'otesis, no una repetici\'on de la anterior.

\subsection{Camino 3 --- M\'as partidas o una partici\'on de validaci\'on}

Con 20 partidas, el IC~95\,\% de cualquier $\Delta R^2$ tiene un ancho
de $\sim$0.13--0.15 (Secci\'on~\ref{sec:bootstrap}) --- suficiente para
que $H_{OPT\_B}$ se distinga de $H_{M1}$/$H_{OPT\_A}$ con traslape
angosto, pero \textbf{no} para resolver del todo el empate persistente
entre estos \'ultimos dos. M\'as partidas (o separar expl\'icitamente
un conjunto de validaci\'on nunca visto durante la b\'usqueda de
coeficientes) es el \'unico camino que ataca directamente el ancho del
intervalo de confianza, en vez de seguir buscando coeficientes.
\textbf{Costo}: depende de cu\'antas partidas nuevas haya disponibles
--- KataGo ya est\'a cacheado como pipeline, as\'i que el costo marginal
es solo el de correr KataGo sobre las partidas nuevas (no hay que
tocar el c\'odigo).

\subsection{Camino 4 --- Completar el cat\'alogo del visor}

Integrar al visor interactivo los 25 candidatos \texttt{sparse\_cubic}
generados al azar (Secci\'on~3.2), pendientes desde el inicio del
experimento. \textbf{Costo}: bajo --- es trabajo de integraci\'on
(estos candidatos ya fueron evaluados, solo falta agregarlos al
cat\'alogo de 48), no c\'omputo nuevo de KataGo ni de
\texttt{relaxation\_field}.

\subsection{Camino 5 --- Categor\'ias 6 y 7: mecanismo nuevo}

Grupos de piedras/vida-muerte y frontera activa/aji
(Secci\'on~\ref{sec:categorias}) son cualitativamente distintos de
moyo/territorio/fase/joseki/fuseki/sente-gote: no se pueden resolver
reutilizando \texttt{relaxation\_field} punto por punto --- necesitan
agregaci\'on a nivel de \textbf{grupo completo de piedras} (vida/muerte
de un grupo no es una propiedad de un punto individual). \textbf{Costo}:
el m\'as alto de los cinco caminos --- requiere dise\~nar un mecanismo
nuevo desde cero, no una extensi\'on del pipeline existente.

\vfill
\noindent{\color{cRule}\hrule height 0.5pt}\vspace{4pt}
{\color{cMuted}\small\itshape
  Pipeline:~\texttt{experiments/07\_moyo\_dataset/}
  \quad\textbullet\quad
  Teor\'ia:~\texttt{experiments/08\_teoria\_invariantes/}
  \quad\textbullet\quad
  20 partidas profesionales \texttt{.sgf}
  \quad\textbullet\quad
  Motor:~KataGo \texttt{kata1-b15c192} (local, OpenCL)
  \quad\textbullet\quad
  Verificaci\'on simb\'olica:~\texttt{sympy}
}

\end{document}
